% language package supplied by unit wrapper

% UTF-8 input supplied by unit wrapper

%\usepackage[utf8]{inputenc}

\usepackage{ifthen}

%\usepackage{fancyhdr}

%Aus Grundvorspann

%Packages

\usepackage{amscd}

\usepackage{amssymb}

\usepackage{epsfig} %Graphik

\usepackage{verbatim,moreverb} %

\usepackage{enumitem}

\usepackage{eurosym}

\usepackage{epigraph}

\setcounter{MaxMatrixCols}{20}

%Mathematische Einzelbefehle

%[[Projekt:Semantische Vorlagen/Mathematischer Modus/Latex]]

\newcommand{\mathdisp}[2]{\[ #1 \, #2 \]}

\newcommand{\mathdisplayteile}[3]{\[ #1 \] \[ #2 \, #3 \]}

\newcommand{\mathind}[3]{$ #1$, $#2 $#3}

\newcommand{\mathbed}[8]{$ #1 ${\ifthenelse{\equal{#2}{}}{, }{ #2 }}$ #3 ${\ifthenelse{\equal{#5}{}}{}{{\ifthenelse{\equal{#4}{}}{, }{ #4 }}}$ #5 $}{\ifthenelse{\equal{#7}{}}{}{{\ifthenelse{\equal{#6}{}}{, }{ #6 }}}$ #7 $}#8}

\newcommand{\mathbeddisp}[8]{\[ #1 {\ifthenelse{\equal{#2}{}}{,\, }{ \text{ #2 } }} #3 {\ifthenelse{\equal{#5}{}}{}{{\ifthenelse{\equal{#4}{}}{,\, }{ \text{ #4 } }}} #5 }{\ifthenelse{\equal{#7}{}}{}{{\ifthenelse{\equal{#6}{}}{,\, }{ \text{ #6 } }}} #7 }#8 \]}

\newcommand{\mathbedtermdisp}[8]{\[ #1\] #2 $ #3 $#4 $ #5 $#6$ #7 $#8}

\newcommand{\mathkon}[4]{$ #1 $ #2 $ #3 $#4}

\newcommand{\mathkor}[5]{{\ifthenelse{\equal{#1}{}}{}{#1 }}$ #2 $ #3 $ #4 $#5}

\newcommand{\mathkork}[5]{{\ifthenelse{\equal{#1}{}}{}{#1 }}$ #2 $ #3 $ #4 $#5}

\newcommand{\mathlist}[5]{$ #1 $\ifthenelse{\equal{#2}{}}{,}{ #2} $ #3 $\ifthenelse{\equal{#4}{}}{,}{ #4} $ #5 $}

\newcommand{\mathlistdisp}[5]{\[  #1  \ifthenelse{\equal{#2}{}}{,}{ \text{ #2 }}  #3  \ifthenelse{\equal{#4}{}}{,}{\text{ #4 } }  #5 \] }

\newcommand{\alignier}[2]{\begin{align*} #1 \, #2 \end{align*}}

\newcommand{\mathlistvierdisp}[7]{\[ #1,\, #3, \, #5 \text{ #6 } #7  \]}

%[[Projekt:Semantische Vorlagen/Mathematische Vergleichsvorlagen/Latex]]

\newcommand{\mavergleichskette}[4]{$ #1 #2 #3#4$}

\newcommand{\mavergleichskettedisp}[4]{\[  #1 #2 #3 #4 \]}

\newcommand{\mavergleichskettedisplang}[4]{\[  #1 #2 #3 #4 \]}

\newcommand{\mavergleichskettealign}[4]{ \begin{eqnarray*}  #1 #2 #3 #4 \end{eqnarray*}  }

\newcommand{\vergleichskette}[9]{ #1 \, #2 \, #3 \ifthenelse
{\equal {#4}{ }} {} {\, #4 \, #5 }  \ifthenelse {\equal {#6}{ } }{}
{\, #6 \, #7 } \ifthenelse {\equal {#8}{ }}{} {\, #8 \, #9 }  }

\newcommand{\vergleichskettelang}[9]{ #1 \, #2 \, #3 \ifthenelse
{\equal {#4}{ }} {} {\, #4 \, #5 }  \ifthenelse {\equal {#6}{ } }{}
{\, #6 \, #7 } \ifthenelse {\equal {#8}{ }}{} {\, #8 \, #9 }  }

\newcommand{\vergleichskettefortsetzung}[8]{ \, #1 \, #2 \ifthenelse
{\equal {#3}{ }} {} {\, #3 \, #4 }  \ifthenelse {\equal {#5}{ } }{}
{\, #5 \, #6 } \ifthenelse {\equal {#7}{ }}{} {\, #7 \, #8 } }

\newcommand{\vergleichskettealign}[9]{ #1 & #2 & #3 \ifthenelse
{\equal {#4}{ }} {} {\cr & #4 & #5 }  \ifthenelse {\equal {#6}{ }
}{} {\cr & #6 & #7 } \ifthenelse {\equal {#8}{ }}{} {\cr & #8 & #9 }
}

\newcommand{\vergleichskettefortsetzungalign}[8]{ \cr & #1 & #2 \ifthenelse
{\equal {#3}{ }} {} {\cr & #3 & #4 }  \ifthenelse {\equal {#5}{ }
}{} {\cr & #5 & #6 } \ifthenelse {\equal {#7}{ }}{} {\cr & #7 & #8 }
}

\newcommand{\mavergleichskettealigndrucklinks}[4]{ \begin{eqnarray*}  #1 #2 #3 #4 \end{eqnarray*}  }

\newcommand{\vergleichskettealigndrucklinks}[9]{  &  & #1 \cr & #2 & #3 \ifthenelse
{\equal {#4}{ }} {} {\cr & #4 & #5 }  \ifthenelse {\equal {#6}{ }
}{} {\cr & #6 & #7 } \ifthenelse {\equal {#8}{ }}{} {\cr & #8 & #9 }
}

\newcommand{\mavergleichskettealignhandlinks}[4]{ \begin{eqnarray*}  #1 #2 #3 #4 \end{eqnarray*}  }

\newcommand{\vergleichskettealignhandlinks}[9]{ #1 & #2 & #3 \ifthenelse
{\equal {#4}{ }} {} {\cr & #4 & #5 }  \ifthenelse {\equal {#6}{ }
}{} {\cr & #6 & #7 } \ifthenelse {\equal {#8}{ }}{} {\cr & #8 & #9 }
}

\newcommand{\mavergleichskettedisphandlinks}[4]{ \[ #1 #2 #3 #4 \] }

\newcommand{\vergleichskettedisphandlinks}[9]{  #1 \, #2 \, #3 \ifthenelse
{\equal {#4}{ }} {} {\, #4 \, #5 }  \ifthenelse {\equal {#6}{ } }{}
{\, #6 \, #7 } \ifthenelse {\equal {#8}{ }}{} {\, #8 \, #9 }  }

\newcommand{\zeilemitteil}[2]{ \ifthenelse {\equal {#2} {} } {#1} { #1 \cr & & #2}   }

\newcommand{\mavergleichskettek}[4]{$ #1 #2 #3#4$}

\newcommand{\vergleichskettek}[9]{ #1 \, #2 \, #3 \ifthenelse
{\equal {#4}{ }} {} {\, #4 \, #5 }  \ifthenelse {\equal {#6}{ } }{}
{\, #6 \, #7 } \ifthenelse {\equal {#8}{ }}{} {\, #8 \, #9 }  }

\newcommand{\vergleichskettefortsetzungk}[8]{ \, #1 \, #2 \ifthenelse
{\equal {#3}{ }} {} {\, #3 \, #4 }  \ifthenelse {\equal {#5}{ } }{}
{\, #5 \, #6 } \ifthenelse {\equal {#7}{ }}{} {\, #7 \, #8 } }

%[[Projekt:Semantische Vorlagen/Mathematische Strukturvorlagen/Latex]]

\newcommand{\N}   {{\mathbb N}}
\newcommand{\Z}   {{\mathbb Z}}
\newcommand{\Q}   {{\mathbb Q}}
\newcommand{\R}   {{\mathbb R}}
\newcommand{\C}   {{\mathbb C}}
\newcommand{\Complex}  {{\mathbb C}}

\newcommand{\defeq}{:=}

\newcommand{\defeqr}{=:}

\newcommand{\maabb}[4]{${\ifthenelse {\equal {#1}{} } {} {#1 \colon }}\,  #2 \rightarrow #3 #4 $}

\newcommand{\maabbele}[6]{$ {\ifthenelse {\equal {#1}{} } {} {#1 \colon }}\,   #2 \rightarrow  #3 , \, #4 \mapsto #5 #6 $}

\newcommand{\maabbdisp}[4]{\[ {\ifthenelse {\equal {#1}{} } {} {#1 \colon }}\,  #2 \longrightarrow #3 #4 \]}

\newcommand{\maabbeledisp}[6]{\[ {\ifthenelse {\equal {#1}{} } {} {#1 \colon }}\,   #2 \longrightarrow  #3 , \, #4 \longmapsto #5 #6 \]}

%Variante bei ifequal-Problem

\newcommand{\maabbnamedisp}[4]{\[ #1 \colon \, #2 \longrightarrow #3 #4 \]}

\newcommand{\maabbeledispvar}[6]{\[ {\ifthenelse {\equal {#1}{} } {} {#1 \colon }}\,   #2 \longrightarrow  #3 , \, #4 \longmapsto #5 #6  \]}

\newcommand{\maabbelementzeiledisplay}[6]{\[ {\ifthenelse {\equal {#1}{} } {} {#1 \colon }}\, #2 \longrightarrow #3 , \]  	\[ #4 \longmapsto #5 #6 \]}

\newcommand{\maabbelementdoppelzeiledisplay}[6]{\[ {\ifthenelse {\equal {#1}{} } {} {#1 \colon }}\, #2 \longrightarrow #3 , \] \[ #4 \longmapsto \] \[  #5 #6 \]}

\newcommand{ \arccot  } {\operatorname{arccot } }

\newcommand{\betrag}[1]{\left| #1 \right|}

\newcommand{\strukturgarbe}{ {\mathcal O} }

\newcommand{\schnittring}[1]{\Gamma (#1 , {\mathcal O})}

\newcommand{\tensor}{\otimes}

%[[Projekt:Semantische Vorlagen/Textgestaltung/Latex]]

\newcommand{\einrueckung}[1]{\par \hspace{1cm}#1 \par}

\newcommand{\zusatz}[2]{(#1)#2}

\newcommand{\zusatzklammer}[3]{(#1#2)#3}

\newcommand{\zusatzgs}[3]{\-- #1 \-- #3}

\newcommand{\zusatzfussnote}[3]{#3\footnote{#1#2}}

%Feinabstimmung

\newcommand{\latexdruckeinskurz}{\!}

\newcommand{\latexdruckzweikurz}{\! \!}

\newcommand{\latexdruckdreikurz}{\! \! \!}

%[[Projekt:Semantische Vorlagen/Satzumgebungen/Latex]]

\newtheorem{fakt}{Fakta}[section] %so wird abschnittsweise %durchnummeriert. Bei [] wird % einfach nummeriert.

\newtheorem{Fakt}{fakt}[Fakt]

\newtheorem{Proposition}[fakt]{Proposisi}

\newtheorem{Korollar}[fakt]{Korolari}

\newtheorem{Lemma}[fakt]{Lema}

\newtheorem{Satz}[fakt]{Teorema}

\newtheorem{Theorem}[fakt]{Teorema}

\theoremstyle{definition}

\newtheorem{Axiom}[fakt]{Aksioma}

\newtheorem{aufgabe}[fakt]{Soal}

\newtheorem{Aufgabe}[fakt]{Soal}

\newtheorem{klausuraufgabe}{Soal ujian}

\newtheorem{Klausuraufgabe}{Soal ujian}

\newtheorem{beispiel}[fakt]{Contoh}

\newtheorem{Beispiel}[fakt]{Contoh}

\newtheorem{Konstruktion}[fakt]{Konstruksi}

\newtheorem{konstruktion}[fakt]{Konstruksi}

\newtheorem{Verfahren}[fakt]{Verfahren}

\newtheorem{verfahren}[fakt]{Verfahren}

\newtheorem{bemerkung}[fakt]{Catatan}

\newtheorem{Bemerkung}[fakt]{Catatan}

\newtheorem{definition}[fakt]{Definisi}

\newtheorem{Definition}[fakt]{Definisi}

\newtheorem{notation}[fakt]{Notation}

\newtheorem{Notation}[fakt]{Notation}

\newtheorem{Frage}[fakt]{Frage}

\newtheorem{frage}[fakt]{Frage}

\newtheorem{problem}[fakt]{Problem}

\newtheorem{Problem}[fakt]{Problem}

\newtheorem{situation}[fakt]{Situation}

\newtheorem{Situation}[fakt]{Situation}

%Einlesegestaltung

%[[Projekt:Semantische Vorlagen/Umgebungsdesign/Latex]]

%Aufgabengestaltung

\newcommand{\inputaufgabe}[4]{\aufgabevorskip \begin{aufgabe} \punkte {#1} \aufgabepunktskip #2 #3 #4\end{aufgabe} \aufgabenachskip}

\newcommand{\inputaufgabegibtloesung}[4]{\aufgabevorskip \begin{aufgabe}\hspace{-1.8 mm}* \punkte {#1} \aufgabepunktskip #2 #3 #4\end{aufgabe} \aufgabenachskip}

\newcommand{\inputaufgabepunktenummervonhand}[3]{\aufgabevorskip {\bf Aufgabe #1} (#2 Punkte) \aufgabepunktskip #3  \aufgabenachskip}

\newcommand{\inputaufgabeloesung}[2]{ \begin{aufgabe} #1 \end{aufgabe} \aufgabenachskip

\bigskip

Lösung

\bigskip #2

}

\newcommand{\inputaufgabeklausurloesung}[3]{\par \newpage \begin{aufgabe} {\ifthenelse {\equal {#1}{}}{} {\ifthenelse {\equal {#1}{1}} {(1 Punkt)} {(#1 Punkte)}}} \par \bigskip #2 \end{aufgabe} \underline{L\"osung:} \par \bigskip #3 }

\newcommand{\inputaufgabeloesungvar}[2]{\aufgabevorskip Aufgabe: #1 \aufgabenachskip
Lösung: #2}

\newcommand{\inputaufgabepunkteloesung}[3]{\aufgabevorskip \begin{aufgabe} {\ifthenelse {\equal {#1}{}}{} {\ifthenelse {\equal {#1}{1}} {(1 Punkt)} {(#1 Punkte)}}}  \aufgabepunktskip #2 \end{aufgabe} \aufgabenachskip \loesung #3}

\newcommand{\loesung}[1]{\underline{Lösung:} #1}

\newcommand{\punkte}[1]{\ifthenelse {\equal {#1}{}}{} {\ifthenelse {\equal {#1}{1}} {(1 poin)} {(#1 poin)}}}

\newcommand{\inputexercise}[4]{\aufgabevorskip \begin{exercise} \points {#1} \aufgabepunktskip #2 #3 #4\end{exercise} \aufgabenachskip}

\newcommand{\points}[1]{\ifthenelse {\equal {#1}{}}{} {\ifthenelse {\equal {#1}{1}} {(1 point)} {(#1 points)}}}

%Beispielgestaltung

\newcommand{\inputbeispiel}[2]
{\beispielvorskip
\begin{beispiel}
\beispielbenennung{#1}
\beispielbenennungnachskip #2
\end{beispiel}
\beispielnachskip}

\newcommand{\beispielbenennung}[1]{ \ifthenelse {\equal {#1}{} \OR \equal {#1}{   }  }{}{(#1)}     }

%Bemerkunggestaltung

\newcommand{\inputbemerkung}[2]
{\bemerkungvorskip
\begin{bemerkung}
\bemerkungbenennung{#1}
\bemerkungbenennungnachskip #2
\end{bemerkung}
\bemerkungnachskip}

\newcommand{\bemerkungbenennung}[1]{  \ifthenelse {\equal {#1}{} \OR \equal {#1}{   }  }{}{(#1)}}

\newcommand{\inputverfahren}[2]
{\bemerkungvorskip
\begin{verfahren}
\bemerkungbenennung{#1}
\bemerkungbenennungnachskip #2
\end{verfahren}
\bemerkungnachskip}

\newcommand{\inputkonstruktion}[2]
{\bemerkungvorskip
\begin{konstruktion}
\bemerkungbenennung{#1}
\bemerkungbenennungnachskip #2
\end{konstruktion}
\bemerkungnachskip}

\newcommand{\inputfrage}[2]
{\bemerkungvorskip
\begin{frage}
\problembenennung{#1}
\bemerkungbenennungnachskip #2
\end{frage}
\bemerkungnachskip}

\newcommand{\fragebenennung}[1]{  \ifthenelse {\equal {#1}{} \OR \equal {#1}{   }  }{}{(#1)}}

\newcommand{\inputproblem}[2]
{\bemerkungvorskip
\begin{problem}
\problembenennung{#1}
\bemerkungbenennungnachskip #2
\end{problem}
\bemerkungnachskip}

\newcommand{\problembenennung}[1]{  \ifthenelse {\equal {#1}{} \OR \equal {#1}{   }  }{}{(#1)}}

%Situationsgestaltung

\newcommand{\inputsituation}[2]
{\begin{situation}
\situationbenennung{#1} #2
\end{situation}
}

\newcommand{\situationbenennung}[1]{\ifthenelse {\equal {#1}{} \OR \equal {#1}{  }  }{}{(#1)}}

%Definitiongestaltung

\newcommand{\inputdefinition}[2]
{\definitionvorskip
\begin{definition}
\definitionbenennung{#1}
\definitionbenennungnachskip #2
\end{definition}
\definitionnachskip}

\newcommand{\inputaxiom}[2]
{\definitionvorskip
\begin{Axiom}
\definitionbenennung{#1}
\definitionbenennungnachskip #2
\end{Axiom}
\definitionnachskip}

\newcommand{\definitionbenennung}[1]{\ifthenelse {\equal {#1}{} \OR \equal {#1}{  }  }{}{(#1)}}

\newcommand{\inputnotation}[2]
{\definitionvorskip \begin{notation} \definitionbenennung{#1} \definitionbenennungnachskip #2 \end{notation} \definitionnachskip}

%Fakt- und Beweisgestaltung

\renewcommand{\proofname}{\hspace{0cm}{\it Bukti}}

\newcommand{\beweis}[1]{\begin{proof} #1 \end{proof}}

\newcommand{\inputfakt}[4]
{\faktvorskip
\begin{#2}

%\label{#1}   (Absetzen, sonst  kann das folgende dahinter rutschen)

\faktbenennung{#3}
\faktbenennungnachskip #4
\end{#2}
\faktnachskip
}

\newcommand{\inputfaktbeweis}[5]
{\faktvorskip
\begin{#2}

%\label{#1}   (Absetzen, sonst  kann das folgende dahinter rutschen)

\faktbenennung{#3}
\faktbenennungnachskip #4
\end{#2}
\faktnachskip
\beweis{#5}}

\newcommand{\inputfaktbeweisnichtvorgefuehrt}[5]{\inputfaktbeweis {#1} {#2} {#3} {#4} {Dieser Beweis wurde in der Vorlesung nicht vorgef\"uhrt.} }

\newcommand{\inputfaktbeweistrivial}[4]{\inputfaktbeweis {#1} {#2} {#3} {#4} {Das ist trivial.} }

\newcommand{\inputfaktuebergangbeweis}[6]
{\faktvorskip
\begin{#2}

%\label{#1} (Absetzen, sonst kann das
% folgende dahinter rutschen)

\faktbenennung{#3}
\faktbenennungnachskip #4
\end{#2}
\faktnachskip #5 \faktnachskip
\beweis{#6}}

\newcommand{\inputbeweis}[1]{\beweis{#1} }

\newcommand{\faktbenennung}[1]{\ifthenelse {\equal {#1}{} \OR \equal {#1}{  }  }{}{(#1)}}

%Gestaltung der Faktstruktur

\newcommand{\faktsituation}[1]{\faktsituationskip #1}

\newcommand{\faktvoraussetzung}[1]{\faktvoraussetzungskip #1}

\newcommand{\faktvoraussetzungleer}[1]{ \hspace{-0,15cm} }

\newcommand{\faktvoraussetzungpos}[1]{\faktvoraussetzungskip #1}

\newcommand{\faktuebergang}[1]{\faktuebergangskip #1}

\newcommand{\faktuebergangleer}[1]{\faktuebergangskip \hspace{-0,15cm} }

\newcommand{\faktfolgerung}[1]{\faktfolgerungskip #1}

\newcommand{\faktzusatz}[1]{\faktzusatzskip #1}

%[[Projekt:Semantische Vorlagen/Beweisgestaltung/Latex]]

\newcommand{\teilbeweis}[5]{#1#2#3#4#5}

\newcommand{\fallunterscheidung}[5]{#1 #2\ifthenelse {\equal {#3}{}} {} { #3}\ifthenelse {\equal {#4}{}} {} { #4}\ifthenelse {\equal {#5}{}} {} { #5}}

\newcommand{\fallunterscheidungzwei}[2]{#1 #2}

\newcommand{\fallunterscheidungdrei}[3]{#1 #2 #3 }

\newcommand{\fallunterscheidungvier}[4]{#1 #2 #3 #4 }

\newcommand{\fallunterscheidungfuenf}[5]{#1 #2 #3 #4 #5 }

%[[Projekt:Semantische Vorlagen/Stichwortumsetzung/Latex]]

\newcommand{\betonung}[2]{\emph{#1}#2}

\newcommand{\stichwort}[2]{\emph{#1}#2}

\newcommand{\stichwortpraemath}[3]{$#1$-\emph{#2}#3}

\newcommand{\definitionswort}[2]{\emph{#1}#2}

\newcommand{\definitionswortpraemath}[3]{$#1$-\emph{#2}#3}

\newcommand{\definitionswortteil}[2]{\emph{#1}#2}

\newcommand{\definitionswortenp}[2]{\emph{#1}#2}

\newcommand{\definitionsverweis}[3]{#1#3}

\newcommand{\definitionsverweisanfuehrung}[3]{\anfuehrung {#1} {#3} }

%Anführungszeichen

\newcommand{\anfuehrung}[2]{"`{#1}"'#2}

\newcommand{\anfuehrungenglisch}[2]{``{#1}''#2}

%[[Projekt:Semantische Auflistungsvorlagen/Latex]]

%Aufzählungen

\newcommand{\aufzaehlungeins}[1]{\begin{enumerate} \item  #1 \end{enumerate}}

\newcommand{\aufzaehlungzwei}[2]{\begin{enumerate} \item  #1 \item  #2 \end{enumerate}}

\newcommand{\aufzaehlungdrei}[3]{\begin{enumerate} \item  #1 \item  #2 \item  #3  \end{enumerate}}

\newcommand{\aufzaehlungvier}[4]{\begin{enumerate} \item  #1 \item  #2 \item  #3 \item  #4 \end{enumerate}}

\newcommand{\aufzaehlungfuenf}[5]{\begin{enumerate} \item  #1 \item  #2 \item  #3 \item  #4 \item  #5  \end{enumerate}}

\newcommand{\aufzaehlungsechs}[6]{\begin{enumerate} \item  #1 \item  #2 \item  #3 \item  #4 \item  #5 \item  #6 \end{enumerate}}

\newcommand{\aufzaehlungsieben}[7]{\begin{enumerate} \item  #1 \item  #2 \item  #3 \item  #4 \item  #5 \item #6 \item  #7 \end{enumerate}}

\newcommand{\aufzaehlungacht}[8]{\begin{enumerate} \item  #1 \item  #2 \item  #3 \item  #4 \item  #5 \item #6 \item  #7 \item #8 \end{enumerate}}

\newcommand{\aufzaehlungneun}[9]{\begin{enumerate} \item  #1 \item  #2 \item  #3 \item  #4 \item  #5 \item #6 \item  #7  \item  #8  \item  #9 \end{enumerate}}

\newcommand{\itemfuenf}[5]{\item  #1 \item  #2 \item  #3 \item  #4 \item  #5}

\newcommand{\itemsechs}[6]{\item  #1 \item  #2 \item  #3 \item  #4 \item  #5 \item  #6}

\newcommand{\itemsieben}[7]{\item  #1 \item  #2 \item  #3 \item  #4 \item  #5 \item  #6 \item #7}

\newcommand{\aufzaehlungeinsabc}[1]{\begin{enumerate}[label=(\alph*)]  #1 \end{enumerate}}

\newcommand{\aufzaehlungzweiabc}[2]{\begin{enumerate}[label=(\alph*)] \item  #1 \item  #2 \end{enumerate}}

\newcommand{\aufzaehlungdreiabc}[3]{\begin{enumerate}[label=(\alph*)] \item  #1 \item  #2 \item  #3  \end{enumerate}}

\newcommand{\aufzaehlungvierabc}[4]{\begin{enumerate}[label=(\alph*)] \item  #1 \item  #2 \item  #3 \item  #4 \end{enumerate}}

\newcommand{\aufzaehlungfuenfabc}[5]{\begin{enumerate}[label=(\alph*)] \item  #1 \item  #2 \item  #3 \item  #4 \item  #5  \end{enumerate}}

\newcommand{\aufzaehlungsechsabc}[6]{\begin{enumerate}[label=(\alph*)] \item  #1 \item  #2 \item  #3 \item  #4 \item  #5 \item  #6 \end{enumerate}}

\newcommand{\aufzaehlungsiebenabc}[7]{\begin{enumerate}[label=(\alph*)] \item  #1 \item  #2 \item  #3 \item  #4 \item  #5 \item #6 \item  #7 \end{enumerate}}

\newcommand{\aufzaehlungachtabc}[8]{\begin{enumerate}[label=(\alph*)] \item  #1 \item  #2 \item  #3 \item  #4 \item  #5 \item #6 \item  #7 \item #8 \end{enumerate}}

\newcommand{\aufzaehlungneunabc}[9]{\begin{enumerate}[label=(\alph*)] \item  #1 \item  #2 \item  #3 \item  #4 \item  #5 \item #6 \item  #7  \item  #8  \item  #9 \end{enumerate}}

\newcommand{\aufzaehlungzweireihe}[2]{\begin{enumerate} #1 #2 \end{enumerate}}

\newcommand{\aufzaehlungdreibeweis}[3]{(1) #1 \par (2) #2 \par (3) #3 }

\newcommand{\aufzaehlungvierbeweis}[4]{(1) #1 \par (2) #2 \par (3) #3  \par (4) #4}

\newcommand{\aufzaehlungfuenfbeweis}[5]{(1) #1 \par (2) #2 \par (3) #3  \par (4) #4  \par (5) #5  }

%Auflistungen

\newcommand{\listitem}{\par \listskip \listdot}

\newcommand{\auflistungzwei}[2]{\par \listskip \listdot  #1 \par \listskip \listdot  #2 \par \listskip}

\newcommand{\auflistungdrei}[3]{\par \listskip \listdot  #1 \par \listskip \listdot  #2\par \listskip \listdot  #3\par\listskip }

\newcommand{\auflistungvier}[4]{\par \listskip \listdot  #1 \par \listskip \listdot  #2\par \listskip \listdot  #3\par \listskip \listdot  #4  \par \listskip }

\newcommand{\auflistungfuenf}[5]{\par \listskip \listdot  #1 \par \listskip \listdot  #2\par \listskip \listdot  #3\par \listskip \listdot  #4  \par \listskip  \listdot  #5 \par \listskip   }

\newcommand{\auflistungsechs}[6]{\par \listskip \listdot  #1 \par \listskip \listdot  #2\par \listskip \listdot  #3\par \listskip \listdot  #4  \par \listskip  \listdot  #5 \par \listskip \listdot  #6 \par \listskip}

\newcommand{\auflistungsieben}[7]{\par \listskip \listdot  #1 \par \listskip \listdot  #2\par \listskip \listdot  #3\par \listskip \listdot  #4  \par \listskip  \listdot  #5 \par \listskip \listdot  #6 \par \listskip \listdot  #7 \par \listskip }

\newcommand{\auflistungacht}[8]{\par \listskip \listdot  #1 \par \listskip \listdot  #2\par \listskip \listdot  #3\par \listskip \listdot  #4  \par \listskip  \listdot  #5 \par \listskip \listdot  #6 \par \listskip \listdot  #7 \par \listskip \listdot  #8  \par \listskip}

\newcommand{\auflistungneun}[9]{\par \listskip \listdot  #1 \par \listskip \listdot  #2\par \listskip \listdot  #3\par \listskip \listdot  #4  \par \listskip  \listdot  #5 \par \listskip \listdot  #6 \par \listskip \listdot  #7 \par \listskip \listdot  #8  \par \listskip  \listdot  #9 \par \listskip  }

\newcommand{\listitemfuenf}[5]{\listitem  #1 \listitem  #2 \listitem  #3 \listitem  #4 \listitem  #5}

\newcommand{\auflistungzweireihe}[2]{ #1 #2 }

%Gestaltung der Auflistungen

\newcommand{\listdot}{$\bullet \,$}

%Bildgestaltung

\newcommand{\inputbild}[1]{\includegraphics[scale]{#1}}

%Tabellengestaltung

%[[Projekt:Semantische Vorlagen/Zeilenkonfiguration/Latex]]

\newcommand{\zeileundzwei}[2]{#1 & #2}

\newcommand{\zeileunddrei}[3]{#1 & #2 & #3}

\newcommand{\zeileundvier}[4]{#1 & #2 & #3 & #4}

\newcommand{\zeileundfuenf}[5]{#1 & #2 & #3 & #4 & #5}

\newcommand{\zeileundsechs}[6]{#1 & #2 & #3 & #4 & #5 & #6}

\newcommand{\zeileundsieben}[7]{#1 & #2 & #3 & #4 & #5 & #6 & #7}

\newcommand{\zeileundacht}[8]{#1 & #2 & #3 & #4 & #5 & #6 & #7 & #8}

\newcommand{\zeileundneun}[9]{#1 & #2 & #3 & #4 & #5 & #6 & #7 & #8 & #9}

\newcommand{\mazeileundzwei}[2]{$#1$ & $#2$}

\newcommand{\mazeileunddrei}[3]{$#1$ & $ #2$ & $#3$}

\newcommand{\mazeileundvier}[4]{$#1$ & $#2$ & $#3$ & $#4$}

\newcommand{\mazeileundfuenf}[5]{$#1$ & $#2$ & $#3$ & $#4$ & $#5$}

\newcommand{\mazeileundsechs}[6]{$#1$ & $#2$ & $#3$ & $#4$ & $#5$ & $#6$}

\newcommand{\mazeileundsieben}[7]{$#1$ & $#2$ & $#3$ & $#4$ & $#5$ & $#6$ & $#7$}

\newcommand{\mazeileundacht}[8]{$#1$ & $#2$ & $#3$ & $#4$ & $#5$ & $#6$ & $#7$ & $#8$}

\newcommand{\mazeileundneun}[9]{$#1$ & $#2$ & $#3$ & $#4$ & $#5$ & $#6$ & $#7$ & $#8$ & $#9$}

%besser als Listen?
%Liste mit und (&) bzw. Bruch.

\newcommand {\listesechsmaund}[6]{$#1$ & $#2$& $#3$ & $#4$ & $#5$ & $#6$}

\newcommand {\listesiebenmaund}[7]{$#1$ & $#2$& $#3$ & $#4$ & $#5$ & $#6$ & $#7$}

\newcommand{\listezweiund}[2]{#1}

\newcommand{\listedreiund}[3]{#1}

\newcommand{\listesiebenund}[7]{#1}

\newcommand{\listesechsbruch}[6]{#1 \\ \hline #2 \\ \hline #3 \\ \hline #4 \\ \hline #5 \\ \hline #6}

\newcommand{\leitzeileunddrei}[3]{#1}

\newcommand{\madoppelzeileunddrei}[6]{$#1$&$#2$&$#3$\\ \hline $#4$&$#5$&$#6$}

%[[Projekt:Semantische Vorlagen/Tabellen/Latex]]

\newcommand{\tabelleviervier}[4]{\par \bigskip \begin{tabular}{|c|c|c|c|} \hline #1 \\ \hline #2 \\ \hline #3 \\ \hline #4 \\ \hline \end{tabular}}

\newcommand{\tabellefuenfvier}[5]{\par \bigskip \begin{tabular}{|c|c|c|c|c|} \hline #1 \\ \hline #2 \\ \hline #3 \\ \hline #4 \\ \hline #5  \\ \hline \end{tabular}}

%Mit math. Einträgen - ohne Leitzeile

\newcommand{\matabellezweivier}[5]{\begin{center} \begin{tabular}{|c|c|} \hline #2 \\ \hline #3 \\ \hline #4 \\ \hline #5  \\ \hline \end{tabular} \end{center} }

\newcommand{\matabellezweifuenf}[6]{\begin{center} \begin{tabular}{|c|c|} \hline     #2 \\ \hline #3 \\ \hline #4 \\ \hline #5 \\ \hline #6  \\ \hline \end{tabular} \end{center} }

\newcommand{\matabellezweisechs}[7]{\begin{center} \begin{tabular}{|c|c|} \hline     #2 \\ \hline #3 \\ \hline #4 \\ \hline #5 \\ \hline #6 \\ \hline #7   \\ \hline \end{tabular} \end{center} }

\newcommand{\matabellezweisieben}[8]{\begin{center} \begin{tabular}{|c|c|} \hline     #2 \\ \hline #3 \\ \hline #4 \\ \hline #5 \\ \hline #6 \\ \hline #7 \\ \hline #8  \\ \hline \end{tabular} \end{center} }

%Mit math. Einträgen - Mit Leitzeile

\newcommand{\matabellesiebenzwoelf}[3]{\begin{center} \begin{tabular}{|c||c|c|c|c|c|c|c|}  \hline  #1   \\ \hline \hline #2 \\ \hline #3 \\ \hline \end{tabular} \end{center} }

\newcommand{\matabellezwoelfzwoelf}[3]{\begin{center} \begin{tabular}{|c||c|c|c|c|c|c|c|c|c|c|c|c|}  \hline  #1   \\ \hline \hline #2 \\ \hline #3 \\ \hline \end{tabular} \end{center} }

\newcommand{\matabelledreineun}[5]{\begin{center} \begin{tabular}{|c|c|c|}\hline  #1 \\ \hline  #2 \\ \hline #3 \\ \hline   #4 \\ \hline  #5 \\ \hline \end{tabular} \end{center} }

\newcommand{\matabelledreizehn}[6]{\begin{center} \begin{tabular}{|c|c|c|}\hline  #1 \\ \hline  #2 \\ \hline #3 \\ \hline   #4 \\ \hline  #5 \\ \hline  #6 \\ \hline \end{tabular} \end{center} }

\newcommand{\matabelledreielf}[6]{\begin{center} \begin{tabular}{|c|c|c|}\hline  #1 \\ \hline  #2 \\ \hline #3 \\ \hline   #4 \\ \hline  #5 \\ \hline  #6 \\ \hline \end{tabular} \end{center} }

%[[Projekt:Semantische Vorlagen/Klausur/Punktetabellen/Latex]]

\newcommand{\punktetabellezehn}
{\vspace{3ex} \begin{center} \renewcommand{\arraystretch}{1.5} \tabcolsep=0.78ex \begin{tabular}{@{}|c|c|c|c|c|c|c|c|c|c|c|c|c|}
\hline Aufgabe & 1 & 2 & 3 & 4 & 5 & 6 & 7 & 8 & 9 & 10 & $\sum$
\\ \hline m\"ogliche Pkt. & \aeins & \azwei & \adrei & \avier & \afuenf & \asechs   & \asieben & \aacht  & \aneun & \azehn & \aelf
\\ \hline erhaltene Pkt.  & \phantom{11} & \phantom{11} & \phantom{11} & \phantom{11} & \phantom{11}& \phantom{11} & \phantom{11} & \phantom{11} & \phantom{11} & & \\ \hline \end{tabular} \end{center} \vspace{2ex} }

\newcommand{\punktetabelleelf}
{\vspace{3ex} \begin{center} \renewcommand{\arraystretch}{1.5} \tabcolsep=0.78ex \begin{tabular}{@{}|c|c|c|c|c|c|c|c|c|c|c|c|c|c|}
\hline Aufgabe & 1 & 2 & 3 & 4 & 5 & 6 & 7 & 8 & 9 & 10 & 11 & $\sum$
\\ \hline m\"ogliche Pkt. & \aeins & \azwei & \adrei & \avier & \afuenf & \asechs   & \asieben & \aacht  & \aneun & \azehn & \aelf & \azwoelf
\\ \hline erhaltene Pkt.  & \phantom{11} & \phantom{11} & \phantom{11} & \phantom{11} & \phantom{11}& \phantom{11} & \phantom{11} & \phantom{11} & \phantom{11} & \phantom{11} & & \\ \hline \end{tabular} \end{center} \vspace{2ex} }

\newcommand{\punktetabellezwoelf}
{\vspace{3ex} \begin{center} \renewcommand{\arraystretch}{1.5} \tabcolsep=0.78ex \begin{tabular}{@{}|c|c|c|c|c|c|c|c|c|c|c|c|c|c|c|}
\hline Aufgabe & 1 & 2 & 3 & 4 & 5 & 6 & 7 & 8 & 9 & 10 & 11 & 12    & $\sum$
\\ \hline m\"ogliche Pkt. & \aeins & \azwei & \adrei & \avier & \afuenf & \asechs   & \asieben & \aacht  & \aneun & \azehn & \aelf & \azwoelf & \adreizehn
\\ \hline erhaltene Pkt.  & \phantom{11} & \phantom{11} & \phantom{11} & \phantom{11} & \phantom{11}& \phantom{11} & \phantom{11} & \phantom{11}
& \phantom{11} & \phantom{11} & \phantom{11} & & \\ \hline \end{tabular} \end{center} \vspace{2ex} }

\newcommand{\punktetabelledreizehn}
{\vspace{3ex} \begin{center} \renewcommand{\arraystretch}{1.5} \tabcolsep=0.78ex \begin{tabular}{@{}|c|c|c|c|c|c|c|c|c|c|c|c|c|c|c|c|}
\hline Aufgabe & 1 & 2 & 3 & 4 & 5 & 6 & 7 & 8 & 9 & 10 & 11 & 12 & 13   & $\sum$
\\ \hline m\"ogliche Pkt. & \aeins & \azwei & \adrei & \avier & \afuenf & \asechs   & \asieben & \aacht  & \aneun & \azehn & \aelf & \azwoelf & \adreizehn & \avierzehn
\\ \hline erhaltene Pkt.  & \phantom{11} & \phantom{11} & \phantom{11} & \phantom{11} & \phantom{11}& \phantom{11} & \phantom{11} & \phantom{11} & \phantom{11}
& \phantom{11} & \phantom{11} & \phantom{11} & & \\ \hline \end{tabular} \end{center} \vspace{2ex} }

\newcommand{\punktetabellevierzehn}
{\vspace{3ex} \begin{center} \renewcommand{\arraystretch}{1.5} \tabcolsep=0.78ex \begin{tabular}{@{}|c|c|c|c|c|c|c|c|c|c|c|c|c|c|c|c|c|}
\hline Aufgabe & 1 & 2 & 3 & 4 & 5 & 6 & 7 & 8 & 9 & 10 & 11 & 12 & 13 & 14  & $\sum$
\\ \hline m\"ogliche Pkt. & \aeins & \azwei & \adrei & \avier & \afuenf & \asechs   & \asieben & \aacht  & \aneun & \azehn & \aelf & \azwoelf & \adreizehn & \avierzehn & \afuenfzehn
\\ \hline erhaltene Pkt. & \phantom{11} & \phantom{11} & \phantom{11} & \phantom{11} & \phantom{11} & \phantom{11}& \phantom{11} & \phantom{11} & \phantom{11} & \phantom{11}
& \phantom{11} & \phantom{11} & \phantom{11} & & \\ \hline \end{tabular} \end{center} \vspace{2ex} }

\newcommand{\punktetabellefuenfzehn}
{\vspace{3ex} \begin{center} \renewcommand{\arraystretch}{1.5} \tabcolsep=0.78ex \begin{tabular}{@{}|c|c|c|c|c|c|c|c|c|c|c|c|c|c|c|c|c|c|}
\hline Aufgabe & 1 & 2 & 3 & 4 & 5 & 6 & 7 & 8 & 9 & 10 & 11 & 12 & 13 & 14 & 15 & $\sum$
\\ \hline m\"ogliche Pkt. & \aeins & \azwei & \adrei & \avier & \afuenf & \asechs   & \asieben & \aacht  & \aneun & \azehn & \aelf & \azwoelf & \adreizehn & \avierzehn & \afuenfzehn
& \asechzehn
\\ \hline erhaltene Pkt. & \phantom{11} & \phantom{11} & \phantom{11}
& \phantom{11} & \phantom{11} & \phantom{11} & \phantom{11}& \phantom{11} & \phantom{11} & \phantom{11} & \phantom{11}
& \phantom{11} & \phantom{11} & \phantom{11} & & \\ \hline \end{tabular} \end{center} \vspace{2ex} }

\newcommand{\punktetabellesechzehn}
{\vspace{3ex} \begin{center} \renewcommand{\arraystretch}{1.5} \tabcolsep=0.78ex \begin{tabular}{@{}|c|c|c|c|c|c|c|c|c|c|c|c|c|c|c|c|c|c|}
\hline Aufgabe & 1 & 2 & 3 & 4 & 5 & 6 & 7 & 8 & 9 & 10 & 11 & 12 & 13 & 14 & 15 & 16 & $\sum$
\\ \hline m\"ogliche Pkt. & \aeins & \azwei & \adrei & \avier & \afuenf & \asechs   & \asieben & \aacht  & \aneun & \azehn & \aelf & \azwoelf & \adreizehn & \avierzehn & \afuenfzehn
& \asechzehn   & \asiebzehn
\\ \hline erhaltene Pkt. & \phantom{11} & \phantom{11} & \phantom{11} & \phantom{11}
& \phantom{11} & \phantom{11} & \phantom{11} & \phantom{11}& \phantom{11} & \phantom{11} & \phantom{11} & \phantom{11}
& \phantom{11} & \phantom{11} & \phantom{11} & & \\ \hline \end{tabular} \end{center} \vspace{2ex} }

\newcommand{\punktetabellesiebzehn}
{\vspace{3ex} \begin{center} \renewcommand{\arraystretch}{1.5} \tabcolsep=0.78ex \begin{tabular}{@{}|c|c|c|c|c|c|c|c|c|c|c|c|c|c|c|c|c|c|c|c|}
\hline Aufgabe & 1 & 2 & 3 & 4 & 5 & 6 & 7 & 8 & 9 & 10 & 11 & 12 & 13 & 14 & 15 & 16 & 17 & $\sum$
\\ \hline m\"ogliche Pkt. & \aeins & \azwei & \adrei & \avier & \afuenf & \asechs   & \asieben & \aacht  & \aneun & \azehn & \aelf & \azwoelf & \adreizehn & \avierzehn & \afuenfzehn
& \asechzehn   & \asiebzehn & \aachtzehn
\\ \hline erhaltene Pkt. & \phantom{11} & \phantom{11} & \phantom{11} & \phantom{11}
& \phantom{11} & \phantom{11} & \phantom{11} & \phantom{11} & \phantom{11}& \phantom{11} & \phantom{11} & \phantom{11} & \phantom{11}
& \phantom{11} & \phantom{11} & \phantom{11} & & \\ \hline \end{tabular} \end{center} \vspace{2ex} }

\newcommand{\punktetabelleachtzehn}
{\vspace{3ex} \begin{center} \renewcommand{\arraystretch}{1.5} \tabcolsep=0.78ex \begin{tabular}{@{}|c|c|c|c|c|c|c|c|c|c|c|c|c|c|c|c|c|c|c|c|c|}
\hline Aufgabe & 1 & 2 & 3 & 4 & 5 & 6 & 7 & 8 & 9 & 10 & 11 & 12 & 13 & 14 & 15 & 16 & 17 & 18 & $\sum$
\\ \hline m\"ogliche Pkt. & \aeins & \azwei & \adrei & \avier & \afuenf & \asechs   & \asieben & \aacht  & \aneun & \azehn & \aelf & \azwoelf & \adreizehn & \avierzehn & \afuenfzehn
& \asechzehn   & \asiebzehn & \aachtzehn & \aneunzehn
\\ \hline erhaltene Pkt. & \phantom{11} & \phantom{11} & \phantom{11} & \phantom{11} & \phantom{11}
& \phantom{11} & \phantom{11} & \phantom{11} & \phantom{11} & \phantom{11}& \phantom{11} & \phantom{11} & \phantom{11} & \phantom{11}
& \phantom{11} & \phantom{11} & \phantom{11} & & \\ \hline \end{tabular} \end{center} \vspace{2ex} }

\newcommand{\punktetabelleneunzehn}
{\vspace{3ex} \begin{center} \renewcommand{\arraystretch}{1.5} \tabcolsep=0.78ex \begin{tabular}{@{}|c|c|c|c|c|c|c|c|c|c|c|c|c|c|c|c|c|c|c|c|c|c|}
\hline Aufgabe & 1 & 2 & 3 & 4 & 5 & 6 & 7 & 8 & 9 & 10 & 11 & 12 & 13 & 14 & 15 & 16 & 17 & 18 & 19 & $\sum$
\\ \hline m\"ogliche Pkt. & \aeins & \azwei & \adrei & \avier & \afuenf & \asechs   & \asieben & \aacht  & \aneun & \azehn & \aelf & \azwoelf & \adreizehn & \avierzehn & \afuenfzehn
& \asechzehn   & \asiebzehn & \aachtzehn & \aneunzehn  & \azwanzig
\\ \hline erhaltene Pkt. & \phantom{11} & \phantom{11} & \phantom{11} & \phantom{11} & \phantom{11} & \phantom{11}
& \phantom{11} & \phantom{11} & \phantom{11} & \phantom{11} & \phantom{11}& \phantom{11} & \phantom{11} & \phantom{11} & \phantom{11}
& \phantom{11} & \phantom{11} & \phantom{11} & & \\ \hline \end{tabular} \end{center} \vspace{2ex} }

\newcommand{\punktetabellezwanzig}
{\vspace{3ex} \begin{center} \renewcommand{\arraystretch}{1.5} \tabcolsep=0.78ex \begin{tabular}{@{}|c|c|c|c|c|c|c|c|c|c|c|c|c|c|c|c|c|c|c|c|c|c|c|}
\hline Aufgabe & 1 & 2 & 3 & 4 & 5 & 6 & 7 & 8 & 9 & 10 & 11 & 12 & 13 & 14 & 15 & 16 & 17 & 18 & 19 & 20 & $\sum$
\\ \hline m\"ogliche Pkt. & \aeins & \azwei & \adrei & \avier & \afuenf & \asechs   & \asieben & \aacht  & \aneun & \azehn & \aelf & \azwoelf & \adreizehn & \avierzehn & \afuenfzehn
& \asechzehn   & \asiebzehn & \aachtzehn & \aneunzehn  & \azwanzig & \aeinundzwanzig
\\ \hline erhaltene Pkt. & \phantom{11} & \phantom{11} & \phantom{11} & \phantom{11} & \phantom{11} & \phantom{11} & \phantom{11}
& \phantom{11} & \phantom{11} & \phantom{11} & \phantom{11} & \phantom{11}& \phantom{11} & \phantom{11} & \phantom{11} & \phantom{11}
& \phantom{11} & \phantom{11} & \phantom{11} & & \\ \hline \end{tabular} \end{center} \vspace{2ex} }

\newcommand{\punktetabelleeinundzwanzig}
{\vspace{3ex} \begin{center} \renewcommand{\arraystretch}{1.5} \tabcolsep=0.78ex \begin{tabular}{@{}|c|c|c|c|c|c|c|c|c|c|c|c|c|c|c|c|c|c|c|c|c|c|c|c|}
\hline Aufgabe & 1 & 2 & 3 & 4 & 5 & 6 & 7 & 8 & 9 & 10 & 11 & 12 & 13 & 14 & 15 & 16 & 17 & 18 & 19 & 20 & 21 & $\sum$
\\ \hline m\"ogliche Pkt. & \aeins & \azwei & \adrei & \avier & \afuenf & \asechs   & \asieben & \aacht  & \aneun & \azehn & \aelf & \azwoelf & \adreizehn & \avierzehn & \afuenfzehn
& \asechzehn   & \asiebzehn & \aachtzehn & \aneunzehn  & \azwanzig & \aeinundzwanzig  & \azweiundzwanzig
\\ \hline erhaltene Pkt. & \phantom{11} & \phantom{11} & \phantom{11} & \phantom{11} & \phantom{11} & \phantom{11} & \phantom{11}
& \phantom{11} & \phantom{11} & \phantom{11} & \phantom{11} & \phantom{11}& \phantom{11} & \phantom{11} & \phantom{11} & \phantom{11}
& \phantom{11} & \phantom{11} & \phantom{11} & \phantom{11} & & \\ \hline \end{tabular} \end{center} \vspace{2ex} }

\newcommand{\punktetabellezweiundzwanzig}
{\vspace{3ex} \begin{center} \renewcommand{\arraystretch}{1.5} \tabcolsep=0.78ex \begin{tabular}{@{}|c|c|c|c|c|c|c|c|c|c|c|c|c|c|c|c|c|c|c|c|c|c|c|c|c|}
\hline Aufgabe & 1 & 2 & 3 & 4 & 5 & 6 & 7 & 8 & 9 & 10 & 11 & 12 & 13 & 14 & 15 & 16 & 17 & 18 & 19 & 20 & 21 & 22 & $\sum$
\\ \hline  Punkte & \aeins & \azwei & \adrei & \avier & \afuenf & \asechs   & \asieben & \aacht  & \aneun & \azehn & \aelf & \azwoelf & \adreizehn & \avierzehn & \afuenfzehn
& \asechzehn   & \asiebzehn & \aachtzehn & \aneunzehn  & \azwanzig & \aeinundzwanzig  & \azweiundzwanzig  & \adreiundzwanzig
\\ \hline Erhalten & \phantom{11} & \phantom{11} & \phantom{11} & \phantom{11} & \phantom{11} & \phantom{11} & \phantom{11}
& \phantom{11} & \phantom{11} & \phantom{11} & \phantom{11} & \phantom{11}& \phantom{11} & \phantom{11} & \phantom{11}  & \phantom{11} & \phantom{11}
& \phantom{11} & \phantom{11} & \phantom{11} & \phantom{11} & & \\ \hline \end{tabular} \end{center} \vspace{2ex} }

%[[Projekt:Semantische Vorlagen/Wahrheitstabellen/Latex]]

\newcommand{\wahrheitstabelleeins}[4]{
\begin{center}
\begin{tabular}{|c|c|}
\multicolumn{2}{l}{\textbf{#1}}\\
\hline
$p$ & $#2$\\ \hline
w  &	#3 \\ \hline
f  &	#4\\ \hline
\end{tabular}
\end{center} }

\newcommand{\wahrheitstabelleeinszwei}[7]{
\begin{center}
\begin{tabular}{|c|c|c|}
\multicolumn{3}{l}{\textbf{#1}}\\
\hline
$p$ & $#2$ & $#5$ \\ \hline
w  &	#3 &  #4 \\ \hline
f  &	#6 &  #7 \\ \hline
\end{tabular}
\end{center} }

\newcommand{\wahrheitstabelleeinsdrei}[4]{ \begin{center} \begin{tabular}{|c|c|c|c|} \multicolumn{4}{l}{\textbf{#1}} \\  \hline #2  \\  \hline #3  \\  \hline  #4  \\  \hline
\end{tabular} \end{center} }

%Wahrheitstabellen mit zwei Variablen p und q

\newcommand{\wahrheitstabellezweieins}[6]{ \begin{center} \begin{tabular}{|c|c|c|} \multicolumn{3}{l}{\textbf{#1}} \\  \hline  #2  \\  \hline  #3  \\  \hline  #4  \\  \hline  #5  \\  \hline  #6  \\  \hline \end{tabular} \end{center} }

\newcommand{\wahrheitstabellezweizwei}[6]{ \begin{center} \begin{tabular}{|c|c|c|c|} \multicolumn{4}{l}{\textbf{#1}} \\  \hline #2  \\  \hline #3  \\  \hline #4  \\  \hline #5  \\  \hline #6  \\  \hline\end{tabular} \end{center} }

\newcommand{\wahrheitstabellezweidrei}[6]{ \begin{center} \begin{tabular}{|c|c|c|c|c|} \multicolumn{5}{l}{\textbf{#1}} \\  \hline #2  \\  \hline #3  \\  \hline #4  \\  \hline #5  \\  \hline #6  \\  \hline \end{tabular} \end{center} }

\newcommand{\wahrheitstabellezweivier}[6]{ \begin{center} \begin{tabular}{|c|c|c|c|c|c|} \multicolumn{6}{l}{\textbf{#1}} \\  \hline #2  \\  \hline #3  \\  \hline #4  \\  \hline  #5  \\  \hline  #6  \\  \hline \end{tabular} \end{center} }

\newcommand{\wahrheitstabellezweifuenf}[6]{ \begin{center} \begin{tabular}{|c|c|c|c|c|c|c|} \multicolumn{7}{l}{\textbf{#1}} \\  \hline #2  \\  \hline #3  \\  \hline #4  \\  \hline #5  \\  \hline #6  \\  \hline
\end{tabular} \end{center} }

%Wahrheitstabellen mit drei Variablen p,q,r

\newcommand{\wahrheitstabelledreieins}[6]{ \begin{center} \begin{tabular}{|c|c|c|c||} \multicolumn{4}{l}{\textbf{#1}} \\  \hline #2  \\  \hline #3  \\  \hline #4  \\  \hline #5  \\  \hline #6 \\  \hline
\end{tabular} \end{center} }

%[[Projekt:Semantische Vorlagen/Wertetabellen/Latex]]

%Wertetabellen

\newcommand{\wertetabellevorskip}{\par \bigskip}

\newcommand{\wertetabellenachskip}{\par \bigskip}

\newcommand{\wertetabelleeinsausteilzeilen}[4]{\wertetabellevorskip \begin{tabular}{|c|c|} \hline #1 & #2  \\ \hline  #3 & #4  \\ \hline \end{tabular} \wertetabellenachskip }

\newcommand{\wertetabellezweiausteilzeilen}[4]{\wertetabellevorskip \begin{tabular}{|c|c|c|} \hline #1 & #2  \\ \hline  #3 & #4  \\ \hline \end{tabular} \wertetabellenachskip }

\newcommand{\wertetabelledreiausteilzeilen}[4]{\wertetabellevorskip \begin{tabular}{|c|c|c|c|} \hline #1 & #2  \\ \hline  #3 & #4  \\ \hline \end{tabular} \wertetabellenachskip }

\newcommand{\wertetabellevierausteilzeilen}[4]{\wertetabellevorskip \begin{tabular}{|c|c|c|c|c|} \hline #1 & #2  \\ \hline  #3 & #4  \\ \hline \end{tabular} \wertetabellenachskip }

\newcommand{\wertetabellefuenfausteilzeilen}[4]{\wertetabellevorskip \begin{tabular}{|c|c|c|c|c|c|} \hline #1 & #2  \\ \hline  #3 & #4  \\ \hline \end{tabular} \wertetabellenachskip }

\newcommand{\wertetabellesechsausteilzeilen}[6]{\wertetabellevorskip \begin{tabular}{|c|c|c|c|c|c|c|} \hline #1 & #2 & #3 \\ \hline  #4 & #5 & #6 \\ \hline \end{tabular} \wertetabellenachskip }

\newcommand{\wertetabellesiebenausteilzeilen}[6]{\wertetabellevorskip \begin{tabular}{|c|c|c|c|c|c|c|c|} \hline #1 & #2 & #3 \\ \hline  #4 & #5 & #6 \\ \hline \end{tabular} \wertetabellenachskip }

\newcommand{\wertetabelleachtausteilzeilen}[6]{\wertetabellevorskip \begin{tabular}{|c|c|c|c|c|c|c|c|c|} \hline #1 & #2 & #3 \\ \hline  #4 & #5 & #6 \\ \hline \end{tabular} \wertetabellenachskip }

\newcommand{\wertetabelleneunausteilzeilen}[6]{\wertetabellevorskip \begin{tabular}{|c|c|c|c|c|c|c|c|c|c|} \hline #1 & #2 & #3 \\ \hline  #4 & #5 & #6 \\ \hline \end{tabular} \wertetabellenachskip }

\newcommand{\wertetabellezehnausteilzeilen}[6]{\wertetabellevorskip \begin{tabular}{|c|c|c|c|c|c|c|c|c|c|c|} \hline #1 & #2 & #3 \\ \hline  #4 & #5 & #6 \\ \hline \end{tabular} \wertetabellenachskip }

\newcommand{\wertetabelleelfausteilzeilen}[8]{\wertetabellevorskip \begin{tabular}{|c|c|c|c|c|c|c|c|c|c|c|c|} \hline #1 & #2 & #3 & #4 \\ \hline #5 & #6 & #7 & #8 \\ \hline \end{tabular} \wertetabellenachskip }

%[[Projekt:Semantische Vorlagen/Verknüpfungstabellen/Latex]]

%Verknüpfungstabellen

\newcommand{\verknuepfungstabelleeins}[2]{\begin{center}
\begin{tabular}{|c||c|}
\hline #1 \\  \hline \hline #2 \\
\hline
\end{tabular}
\end{center} }

\newcommand{\verknuepfungstabellezwei}[3]{\begin{center}
\begin{tabular}{|c||c|c|}
\hline #1 \\  \hline \hline #2 \\ \hline #3 \\
\hline
\end{tabular}
\end{center} }

\newcommand{\verknuepfungstabelledrei}[4]{\begin{center}
\begin{tabular}{|c||c|c|c|}
\hline #1 \\  \hline \hline #2 \\ \hline #3 \\ \hline #4 \\
\hline
\end{tabular}
\end{center} }

\newcommand{\verknuepfungstabellevier}[5]{\begin{center}
\begin{tabular}{|c||c|c|c|c|}
\hline #1 \\  \hline \hline #2 \\ \hline #3 \\ \hline #4 \\  \hline  #5 \\
\hline
\end{tabular}
\end{center} }

\newcommand{\verknuepfungstabellefuenf}[6]{\begin{center}
\begin{tabular}{|c||c|c|c|c|c|}
\hline #1 \\  \hline \hline #2 \\ \hline #3 \\ \hline #4 \\  \hline  #5 \\ \hline  #6 \\
\hline
\end{tabular}
\end{center} }

\newcommand{\verknuepfungstabellesechs}[7]{\begin{center}
\begin{tabular}{|c||c|c|c|c|c|c|}
\hline #1 \\  \hline \hline #2 \\ \hline #3 \\ \hline #4 \\  \hline  #5 \\ \hline  #6 \\  \hline #7 \\
\hline
\end{tabular}
\end{center} }

\newcommand{\verknuepfungstabellesieben}[8]{\begin{center}
\begin{tabular}{|c||c|c|c|c|c|c|c|c|}
\hline #1 \\  \hline \hline #2 \\ \hline #3 \\ \hline #4 \\  \hline  #5 \\ \hline  #6 \\  \hline #7 \\  \hline  #8 \\
\hline
\end{tabular}
\end{center} }

\newcommand{\verknuepfungstabelleacht}[9]{\begin{center}
\begin{tabular}{|c||c|c|c|c|c|c|c|c|}
\hline #1 \\  \hline \hline #2 \\ \hline #3 \\ \hline #4 \\  \hline  #5 \\ \hline  #6 \\  \hline #7 \\  \hline  #8 \\ \hline  #9 \\
\hline
\end{tabular}
\end{center} }

%[[Projekt:Semantische Vorlagen/Initialisierung für Daten/Latex]]

\newcommand{\aeins}{    }

\newcommand{\azwei}{     }

\newcommand{\adrei}{    }

\newcommand{\avier}{    }

\newcommand{\afuenf}{    }

\newcommand{\asechs}{    }

\newcommand{\asieben}{    }

\newcommand{\aacht}{    }

\newcommand{\aneun}{    }

\newcommand{\azehn}{    }

\newcommand{\aelf}{    }

\newcommand{\azwoelf}{    }

\newcommand{\adreizehn}{     }

\newcommand{\avierzehn}{    }

\newcommand{\afuenfzehn}{    }

\newcommand{\asechzehn}{    }

\newcommand{\asiebzehn}{    }

\newcommand{\aachtzehn}{    }

\newcommand{\aneunzehn}{    }

\newcommand{\azwanzig}{    }

\newcommand{\aeinundzwanzig}{    }

\newcommand{\azweiundzwanzig}{    }

\newcommand{\adreiundzwanzig}{    }

\newcommand{\avierundzwanzig}{    }

\newcommand{\afuenfundzwanzig}{    }

\newcommand{\asechsundzwanzig}{    }

\newcommand{\asiebenundzwanzig}{    }

\newcommand{\aachtundzwanzig}{    }

\newcommand{\aneunundzwanzig}{    }

\newcommand{\leitzeilenull}{}

\newcommand{\leitzeileeins}{}

\newcommand{\leitzeilezwei}{}

\newcommand{\leitzeiledrei}{}

\newcommand{\leitzeilevier}{}

\newcommand{\leitzeilefuenf}{}

\newcommand{\leitzeilesechs}{}

\newcommand{\leitzeilesieben}{}

\newcommand{\leitzeileacht}{}

\newcommand{\leitzeileneun}{}

\newcommand{\leitzeilezehn}{}

\newcommand{\leitzeileelf}{}

\newcommand{\leitzeilezwoelf}{}

\newcommand{\leitspaltenull}{}

\newcommand{\leitspalteeins}{}

\newcommand{\leitspaltezwei}{}

\newcommand{\leitspaltedrei}{}

\newcommand{\leitspaltevier}{}

\newcommand{\leitspaltefuenf}{}

\newcommand{\leitspaltesechs}{}

\newcommand{\leitspaltesieben}{}

\newcommand{\leitspalteacht}{}

\newcommand{\leitspalteneun}{}

\newcommand{\leitspaltezehn}{}

\newcommand{\leitspalteelf}{}

\newcommand{\leitspaltezwoelf}{}

\newcommand{\leitspaltedreizehn}{}

\newcommand{\leitspaltevierzehn}{}

\newcommand{\leitspaltefuenfzehn}{}

\newcommand{\leitspaltesechzehn}{}

\newcommand{\leitspaltesiebzehn}{}

\newcommand{\leitspalteachtzehn}{}

\newcommand{\leitspalteneunzehn}{}

\newcommand{\leitspaltezwanzig}{}

\newcommand{\aeinsxeins}{}

\newcommand{\aeinsxzwei}{}

\newcommand{\aeinsxdrei}{}

\newcommand{\aeinsxvier}{}

\newcommand{\aeinsxfuenf}{}

\newcommand{\aeinsxsechs}{}

\newcommand{\aeinsxsieben}{}

\newcommand{\aeinsxacht}{}

\newcommand{\aeinsxneun}{}

\newcommand{\aeinsxzehn}{}

\newcommand{\aeinsxelf}{}

\newcommand{\aeinsxzwoelf}{}

\newcommand{\azweixeins}{}

\newcommand{\azweixzwei}{}

\newcommand{\azweixdrei}{}

\newcommand{\azweixvier}{}

\newcommand{\azweixfuenf}{}

\newcommand{\azweixsechs}{}

\newcommand{\azweixsieben}{}

\newcommand{\azweixacht}{}

\newcommand{\azweixneun}{}

\newcommand{\azweixzehn}{}

\newcommand{\azweixelf}{}

\newcommand{\azweixzwoelf}{}

\newcommand{\adreixeins}{}

\newcommand{\adreixzwei}{}

\newcommand{\adreixdrei}{}

\newcommand{\adreixvier}{}

\newcommand{\adreixfuenf}{}

\newcommand{\adreixsechs}{}

\newcommand{\adreixsieben}{}

\newcommand{\adreixacht}{}

\newcommand{\adreixneun}{}

\newcommand{\adreixzehn}{}

\newcommand{\adreixelf}{}

\newcommand{\adreixzwoelf}{}

\newcommand{\avierxeins}{}

\newcommand{\avierxzwei}{}

\newcommand{\avierxdrei}{}

\newcommand{\avierxvier}{}

\newcommand{\avierxfuenf}{}

\newcommand{\avierxsechs}{}

\newcommand{\avierxsieben}{}

\newcommand{\avierxacht}{}

\newcommand{\avierxneun}{}

\newcommand{\avierxzehn}{}

\newcommand{\avierxelf}{}

\newcommand{\avierxzwoelf}{}

\newcommand{\afuenfxeins}{}

\newcommand{\afuenfxzwei}{}

\newcommand{\afuenfxdrei}{}

\newcommand{\afuenfxvier}{}

\newcommand{\afuenfxfuenf}{}

\newcommand{\afuenfxsechs}{}

\newcommand{\afuenfxsieben}{}

\newcommand{\afuenfxacht}{}

\newcommand{\afuenfxneun}{}

\newcommand{\afuenfxzehn}{}

\newcommand{\afuenfxelf}{}

\newcommand{\afuenfxzwoelf}{}

\newcommand{\asechsxeins}{}

\newcommand{\asechsxzwei}{}

\newcommand{\asechsxdrei}{}

\newcommand{\asechsxvier}{}

\newcommand{\asechsxfuenf}{}

\newcommand{\asechsxsechs}{}

\newcommand{\asechsxsieben}{}

\newcommand{\asechsxacht}{}

\newcommand{\asechsxneun}{}

\newcommand{\asechsxzehn}{}

\newcommand{\asechsxelf}{}

\newcommand{\asechsxzwoelf}{}

\newcommand{\asiebenxeins}{}

\newcommand{\asiebenxzwei}{}

\newcommand{\asiebenxdrei}{}

\newcommand{\asiebenxvier}{}

\newcommand{\asiebenxfuenf}{}

\newcommand{\asiebenxsechs}{}

\newcommand{\asiebenxsieben}{}

\newcommand{\asiebenxacht}{}

\newcommand{\asiebenxneun}{}

\newcommand{\asiebenxzehn}{}

\newcommand{\asiebenxelf}{}

\newcommand{\asiebenxzwoelf}{}

\newcommand{\aachtxeins}{}

\newcommand{\aachtxzwei}{}

\newcommand{\aachtxdrei}{}

\newcommand{\aachtxvier}{}

\newcommand{\aachtxfuenf}{}

\newcommand{\aachtxsechs}{}

\newcommand{\aachtxsieben}{}

\newcommand{\aachtxacht}{}

\newcommand{\aachtxneun}{}

\newcommand{\aachtxzehn}{}

\newcommand{\aachtxelf}{}

\newcommand{\aachtxzwoelf}{}

\newcommand{\aneunxeins}{}

\newcommand{\aneunxzwei}{}

\newcommand{\aneunxdrei}{}

\newcommand{\aneunxvier}{}

\newcommand{\aneunxfuenf}{}

\newcommand{\aneunxsechs}{}

\newcommand{\aneunxsieben}{}

\newcommand{\aneunxacht}{}

\newcommand{\aneunxneun}{}

\newcommand{\aneunxzehn}{}

\newcommand{\aneunxelf}{}

\newcommand{\aneunxzwoelf}{}

\newcommand{\azehnxeins}{}

\newcommand{\azehnxzwei}{}

\newcommand{\azehnxdrei}{}

\newcommand{\azehnxvier}{}

\newcommand{\azehnxfuenf}{}

\newcommand{\azehnxsechs}{}

\newcommand{\azehnxsieben}{}

\newcommand{\azehnxacht}{}

\newcommand{\azehnxneun}{}

\newcommand{\azehnxzehn}{}

\newcommand{\azehnxelf}{}

\newcommand{\azehnxzwoelf}{}

\newcommand{\aelfxeins}{}

\newcommand{\aelfxzwei}{}

\newcommand{\aelfxdrei}{}

\newcommand{\aelfxvier}{}

\newcommand{\aelfxfuenf}{}

\newcommand{\aelfxsechs}{}

\newcommand{\aelfxsieben}{}

\newcommand{\aelfxacht}{}

\newcommand{\aelfxneun}{}

\newcommand{\aelfxzehn}{}

\newcommand{\aelfxelf}{}

\newcommand{\aelfxzwoelf}{}

\newcommand{\azwoelfxeins}{}

\newcommand{\azwoelfxzwei}{}

\newcommand{\azwoelfxdrei}{}

\newcommand{\azwoelfxvier}{}

\newcommand{\azwoelfxfuenf}{}

\newcommand{\azwoelfxsechs}{}

\newcommand{\azwoelfxsieben}{}

\newcommand{\azwoelfxacht}{}

\newcommand{\azwoelfxneun}{}

\newcommand{\azwoelfxzehn}{}

\newcommand{\azwoelfxelf}{}

\newcommand{\azwoelfxzwoelf}{}

\newcommand{\adreizehnxeins}{}

\newcommand{\adreizehnxzwei}{}

\newcommand{\adreizehnxdrei}{}

\newcommand{\adreizehnxvier}{}

\newcommand{\adreizehnxfuenf}{}

\newcommand{\adreizehnxsechs}{}

\newcommand{\adreizehnxsieben}{}

\newcommand{\adreizehnxacht}{}

\newcommand{\adreizehnxneun}{}

\newcommand{\adreizehnxzehn}{}

\newcommand{\adreizehnxelf}{}

\newcommand{\adreizehnxzwoelf}{}

\newcommand{\avierzehnxeins}{}

\newcommand{\avierzehnxzwei}{}

\newcommand{\avierzehnxdrei}{}

\newcommand{\avierzehnxvier}{}

\newcommand{\avierzehnxfuenf}{}

\newcommand{\avierzehnxsechs}{}

\newcommand{\avierzehnxsieben}{}

\newcommand{\avierzehnxacht}{}

\newcommand{\avierzehnxneun}{}

\newcommand{\avierzehnxzehn}{}

\newcommand{\avierzehnxelf}{}

\newcommand{\avierzehnxzwoelf}{}

\newcommand{\afuenfzehnxeins}{}

\newcommand{\afuenfzehnxzwei}{}

\newcommand{\afuenfzehnxdrei}{}

\newcommand{\afuenfzehnxvier}{}

\newcommand{\afuenfzehnxfuenf}{}

\newcommand{\afuenfzehnxsechs}{}

\newcommand{\afuenfzehnxsieben}{}

\newcommand{\afuenfzehnxacht}{}

\newcommand{\afuenfzehnxneun}{}

\newcommand{\afuenfzehnxzehn}{}

\newcommand{\afuenfzehnxelf}{}

\newcommand{\afuenfzehnxzwoelf}{}

\newcommand{\asechzehnxeins}{}

\newcommand{\asechzehnxzwei}{}

\newcommand{\asechzehnxdrei}{}

\newcommand{\asechzehnxvier}{}

\newcommand{\asechzehnxfuenf}{}

\newcommand{\asechzehnxsechs}{}

\newcommand{\asechzehnxsieben}{}

\newcommand{\asechzehnxacht}{}

\newcommand{\asechzehnxneun}{}

\newcommand{\asechzehnxzehn}{}

\newcommand{\asechzehnxelf}{}

\newcommand{\asechzehnxzwoelf}{}

\newcommand{\asiebzehnxeins}{}

\newcommand{\asiebzehnxzwei}{}

\newcommand{\asiebzehnxdrei}{}

\newcommand{\asiebzehnxvier}{}

\newcommand{\asiebzehnxfuenf}{}

\newcommand{\asiebzehnxsechs}{}

\newcommand{\asiebzehnxsieben}{}

\newcommand{\asiebzehnxacht}{}

\newcommand{\asiebzehnxneun}{}

\newcommand{\asiebzehnxzehn}{}

\newcommand{\asiebzehnxelf}{}

\newcommand{\asiebzehnxzwoelf}{}

\newcommand{\aachtzehnxeins}{}

\newcommand{\aachtzehnxzwei}{}

\newcommand{\aachtzehnxdrei}{}

\newcommand{\aachtzehnxvier}{}

\newcommand{\aachtzehnxfuenf}{}

\newcommand{\aachtzehnxsechs}{}

\newcommand{\aachtzehnxsieben}{}

\newcommand{\aachtzehnxacht}{}

\newcommand{\aachtzehnxneun}{}

\newcommand{\aachtzehnxzehn}{}

\newcommand{\aachtzehnxelf}{}

\newcommand{\aachtzehnxzwoelf}{}

%[[Projekt:Semantische Vorlagen/Tabellen mit benannten Parametern/Latex]]

%Tabellen mit zwei Zeilen

\newcommand{\tabelleleitzweixzwei}{

\begin{center}

\begin{tabular}{|c||c|c|}

\hline  \leitzeilenull  &  \leitzeileeins  &  \leitzeilezwei   \\

\hline \hline  \leitspalteeins  &  $ \aeinsxeins $  &  $ \aeinsxzwei $ \\

\hline  \leitspaltezwei  &  $ \azweixeins $  &  $ \azweixzwei $   \\

\hline

\end{tabular}

\end{center} }

%Tabellen mit drei Zeilen

\newcommand{\tabelleleitdreixdrei}{

\begin{center}

\begin{tabular}{|c||c|c|c||}

\hline  \leitzeilenull  &  \leitzeileeins  &  \leitzeilezwei  &  \leitzeiledrei        \\

\hline \hline  \leitspalteeins  &  $ \aeinsxeins $  &  $ \aeinsxzwei $  &  $ \aeinsxdrei $      \\

\hline  \leitspaltezwei  &  $ \azweixeins $  &  $ \azweixzwei $  &  $ \azweixdrei $       \\

\hline  \leitspaltedrei  &  $ \adreixeins $  &  $ \adreixzwei $  &  $ \adreixdrei $       \\

\hline

\end{tabular}

\end{center} }

\newcommand{\tabelleleitdreixvier}{

\begin{center}

\begin{tabular}{|c||c|c|c|c||}

\hline  \leitzeilenull  &  \leitzeileeins  &  \leitzeilezwei  &  \leitzeiledrei  &  \leitzeilevier      \\

\hline \hline  \leitspalteeins  &  $ \aeinsxeins $  &  $ \aeinsxzwei $  &  $ \aeinsxdrei $  &  $ \aeinsxvier $   \\

\hline  \leitspaltezwei  &  $ \azweixeins $  &  $ \azweixzwei $  &  $ \azweixdrei $  &  $ \azweixvier $     \\

\hline  \leitspaltedrei  &  $ \adreixeins $  &  $ \adreixzwei $  &  $ \adreixdrei $  &  $ \adreixvier $     \\

\hline

\end{tabular}

\end{center} }

\newcommand{\tabelleleitdreixsechs}{

\begin{center}

\begin{tabular}{|c||c|c|c|c|c|c|}

\hline  \leitzeilenull  &  \leitzeileeins  &  \leitzeilezwei  &  \leitzeiledrei  &  \leitzeilevier  &  \leitzeilefuenf   &  \leitzeilesechs    \\

\hline \hline  \leitspalteeins  &  $ \aeinsxeins $  &  $ \aeinsxzwei $  &  $ \aeinsxdrei $  &  $ \aeinsxvier $  &  $ \aeinsxfuenf $  &  $ \aeinsxsechs $  \\

\hline  \leitspaltezwei  &  $ \azweixeins $  &  $ \azweixzwei $  &  $ \azweixdrei $  &  $ \azweixvier $  &  $ \azweixfuenf $   &  $ \azweixsechs $  \\

\hline  \leitspaltedrei  &  $ \adreixeins $  &  $ \adreixzwei $  &  $ \adreixdrei $  &  $ \adreixvier $  &  $ \adreixfuenf $   &  $ \adreixsechs $   \\

\hline

\end{tabular}

\end{center} }

\newcommand{\tabelleleitdreixelf}{

\begin{center}

\begin{tabular}{|c||c|c|c|c|c|c|c|c|c|c|c|}

\hline  \leitzeilenull  &  \leitzeileeins  &  \leitzeilezwei  &  \leitzeiledrei  &  \leitzeilevier  &  \leitzeilefuenf   &  \leitzeilesechs  &  \leitzeilesieben  &  \leitzeileacht  &  \leitzeileneun  &  \leitzeilezehn &  \leitzeileelf    \\

\hline \hline  \leitspalteeins  &  $ \aeinsxeins $  &  $ \aeinsxzwei $  &  $ \aeinsxdrei $  &  $ \aeinsxvier $  &  $ \aeinsxfuenf $  &  $ \aeinsxsechs $  &  $ \aeinsxsieben $  &  $ \aeinsxacht $  &  $ \aeinsxneun $ &  $ \aeinsxzehn $  &  $ \aeinsxelf $  \\

\hline  \leitspaltezwei  &  $ \azweixeins $  &  $ \azweixzwei $  &  $ \azweixdrei $  &  $ \azweixvier $  &  $ \azweixfuenf $   &  $ \azweixsechs $  &  $ \azweixsieben $  &  $ \azweixacht $  &  $ \azweixneun $  &  $ \azweixzehn $ &  $ \azweixelf $     \\

\hline  \leitspaltedrei  &  $ \adreixeins $  &  $ \adreixzwei $  &  $ \adreixdrei $  &  $ \adreixvier $  &  $ \adreixfuenf $   &  $ \adreixsechs $  &  $ \adreixsieben $  &  $ \adreixacht $  &  $ \adreixneun $  &  $ \adreixzehn $     &  $ \adreixelf $   \\

\hline

\end{tabular}

\end{center} }

%Tabellen mit vier Zeilen

\newcommand{\tabelleleitvierxzwei}{

\begin{center}

\begin{tabular}{|c||c|c|}

\hline  \leitzeilenull  &  \leitzeileeins  &  \leitzeilezwei   \\

\hline \hline  \leitspalteeins  &  $ \aeinsxeins $  &  $ \aeinsxzwei $    \\

\hline  \leitspaltezwei  &  $ \azweixeins $  &  $ \azweixzwei $    \\

\hline  \leitspaltedrei  &  $ \adreixeins $  &  $ \adreixzwei $    \\

\hline  \leitspaltevier  &  $ \avierxeins $  &  $ \avierxzwei $    \\

\hline

\end{tabular}

\end{center} }

\newcommand{\tabelleleitvierxdrei}{

\begin{center}

\begin{tabular}{|c||c|c|c|}

\hline  \leitzeilenull  &  \leitzeileeins  &  \leitzeilezwei  &  \leitzeiledrei     \\

\hline \hline  \leitspalteeins  &  $ \aeinsxeins $  &  $ \aeinsxzwei $  &  $ \aeinsxdrei $     \\

\hline  \leitspaltezwei  &  $ \azweixeins $  &  $ \azweixzwei $  &  $ \azweixdrei $      \\

\hline  \leitspaltedrei  &  $ \adreixeins $  &  $ \adreixzwei $  &  $ \adreixdrei $    \\

\hline  \leitspaltevier  &  $ \avierxeins $  &  $ \avierxzwei $  &  $ \avierxdrei $    \\

\hline

\end{tabular}

\end{center} }

\newcommand{\tabelleleitvierxvier}{

\begin{center}

\begin{tabular}{|c||c|c|c|c|}

\hline  \leitzeilenull  &  \leitzeileeins  &  \leitzeilezwei  &  \leitzeiledrei  &  \leitzeilevier    \\

\hline \hline  \leitspalteeins  &  $ \aeinsxeins $  &  $ \aeinsxzwei $  &  $ \aeinsxdrei $  &  $ \aeinsxvier $   \\

\hline  \leitspaltezwei  &  $ \azweixeins $  &  $ \azweixzwei $  &  $ \azweixdrei $  &  $ \azweixvier $     \\

\hline  \leitspaltedrei  &  $ \adreixeins $  &  $ \adreixzwei $  &  $ \adreixdrei $  &  $ \adreixvier $   \\

\hline  \leitspaltevier  &  $ \avierxeins $  &  $ \avierxzwei $  &  $ \avierxdrei $  &  $ \avierxvier $   \\

\hline

\end{tabular}

\end{center} }

\newcommand{\tabelleleittextvierxzwei}{

	\begin{center}

		\begin{tabular}{|c||c|c|}

			\hline  \leitzeilenull  &  \leitzeileeins  &  \leitzeilezwei   \\

			\hline \hline  \leitspalteeins  &   \aeinsxeins   &   \aeinsxzwei    \\

			\hline  \leitspaltezwei  &  \azweixeins   &   \azweixzwei     \\

			\hline  \leitspaltedrei  &  \adreixeins   &   \adreixzwei     \\

			\hline  \leitspaltevier  &  \avierxeins   &   \avierxzwei     \\

			\hline

		\end{tabular}

	\end{center} }

\newcommand{\tabelleleitvierxacht}{

\begin{center}

\begin{tabular}{|c||c|c|c|c|c|c|c|c|}

\hline  \leitzeilenull  &  \leitzeileeins  &  \leitzeilezwei  &  \leitzeiledrei  &  \leitzeilevier  &  \leitzeilefuenf  &  \leitzeilesechs  &  \leitzeilesieben  &  \leitzeileacht  \\

\hline \hline  \leitspalteeins  &  $ \aeinsxeins $  &  $ \aeinsxzwei $  &  $ \aeinsxdrei $  &  $ \aeinsxvier $  &  $ \aeinsxfuenf $  &  $ \aeinsxsechs $  &  $ \aeinsxsieben $  &  $ \aeinsxacht $  \\

\hline  \leitspaltezwei  &  $ \azweixeins $  &  $ \azweixzwei $  &  $ \azweixdrei $  &  $ \azweixvier $   &  $ \azweixfuenf $  &  $ \azweixsechs $  &  $ \azweixsieben $  &  $ \azweixacht $   \\

\hline  \leitspaltedrei  &  $ \adreixeins $  &  $ \adreixzwei $  &  $ \adreixdrei $  &  $ \adreixvier $ &  $ \adreixfuenf $  &  $ \adreixsechs $  &  $ \adreixsieben $  &  $ \adreixacht $   \\

\hline  \leitspaltevier  &  $ \avierxeins $  &  $ \avierxzwei $  &  $ \avierxdrei $  &  $ \avierxvier $  &  $ \avierxfuenf $  &  $ \avierxsechs $  &  $ \avierxsieben $  &  $ \avierxacht $  \\

\hline

\end{tabular}

\end{center} }

%Tabellen mit fünf Zeilen

\newcommand{\tabelleleitfuenfxfuenf}{

\begin{center}

\begin{tabular}{|c||c|c|c|c|c|}

\hline  \leitzeilenull  &  \leitzeileeins  &  \leitzeilezwei  &  \leitzeiledrei  &  \leitzeilevier  &  \leitzeilefuenf   \\

\hline \hline  \leitspalteeins  &  $ \aeinsxeins $  &  $ \aeinsxzwei $  &  $ \aeinsxdrei $  &  $ \aeinsxvier $  &  $ \aeinsxfuenf $   \\

\hline  \leitspaltezwei  &  $ \azweixeins $  &  $ \azweixzwei $  &  $ \azweixdrei $  &  $ \azweixvier $  &  $ \azweixfuenf $     \\

\hline  \leitspaltedrei  &  $ \adreixeins $  &  $ \adreixzwei $  &  $ \adreixdrei $  &  $ \adreixvier $  &  $ \adreixfuenf $      \\

\hline  \leitspaltevier  &  $ \avierxeins $  &  $ \avierxzwei $  &  $ \avierxdrei $  &  $ \avierxvier $  &  $ \avierxfuenf $     \\

\hline  \leitspaltefuenf  &  $ \afuenfxeins $  &  $ \afuenfxzwei $  &  $ \afuenfxdrei $  &  $ \afuenfxvier $  &  $ \afuenfxfuenf $    \\

\hline

\end{tabular}

\end{center} }

%Tabellen mit sechs Zeilen

\newcommand{\tabelleleitsechsxdrei}{

\begin{center}

\begin{tabular}{|c||c|c|c|}

\hline  \leitzeilenull  &  \leitzeileeins  &  \leitzeilezwei  &  \leitzeiledrei    \\

\hline \hline  \leitspalteeins  &  $ \aeinsxeins $  &  $ \aeinsxzwei $  &  $ \aeinsxdrei $    \\

\hline  \leitspaltezwei  &  $ \azweixeins $  &  $ \azweixzwei $  &  $ \azweixdrei $      \\

\hline  \leitspaltedrei  &  $ \adreixeins $  &  $ \adreixzwei $  &  $ \adreixdrei $     \\

\hline  \leitspaltevier  &  $ \avierxeins $  &  $ \avierxzwei $  &  $ \avierxdrei $      \\

\hline  \leitspaltefuenf  &  $ \afuenfxeins $  &  $ \afuenfxzwei $  &  $ \afuenfxdrei $     \\

\hline  \leitspaltesechs  &  $ \asechsxeins $  &  $ \asechsxzwei $  &  $ \asechsxdrei $      \\

\hline

\end{tabular}

\end{center} }

\newcommand{\tabelleleitsechsxfuenf}{

\begin{center}

\begin{tabular}{|c||c|c|c|c|c|}

\hline  \leitzeilenull  &  \leitzeileeins  &  \leitzeilezwei  &  \leitzeiledrei  &  \leitzeilevier  &  \leitzeilefuenf   \\

\hline \hline  \leitspalteeins  &  $ \aeinsxeins $  &  $ \aeinsxzwei $  &  $ \aeinsxdrei $  &  $ \aeinsxvier $  &  $ \aeinsxfuenf $   \\

\hline  \leitspaltezwei  &  $ \azweixeins $  &  $ \azweixzwei $  &  $ \azweixdrei $  &  $ \azweixvier $  &  $ \azweixfuenf $     \\

\hline  \leitspaltedrei  &  $ \adreixeins $  &  $ \adreixzwei $  &  $ \adreixdrei $  &  $ \adreixvier $  &  $ \adreixfuenf $      \\

\hline  \leitspaltevier  &  $ \avierxeins $  &  $ \avierxzwei $  &  $ \avierxdrei $  &  $ \avierxvier $  &  $ \avierxfuenf $     \\

\hline  \leitspaltefuenf  &  $ \afuenfxeins $  &  $ \afuenfxzwei $  &  $ \afuenfxdrei $  &  $ \afuenfxvier $  &  $ \afuenfxfuenf $    \\

\hline  \leitspaltesechs  &  $ \asechsxeins $  &  $ \asechsxzwei $  &  $ \asechsxdrei $  &  $ \asechsxvier $  &  $ \asechsxfuenf $    \\

\hline

\end{tabular}

\end{center} }

\newcommand{\tabelleleitsechsxsechs}{

\begin{center}

\begin{tabular}{|c||c|c|c|c|c|c|}

\hline  \leitzeilenull  &  \leitzeileeins  &  \leitzeilezwei  &  \leitzeiledrei  &  \leitzeilevier  &  \leitzeilefuenf  &  \leitzeilesechs   \\

\hline \hline  \leitspalteeins  &  $ \aeinsxeins $  &  $ \aeinsxzwei $  &  $ \aeinsxdrei $  &  $ \aeinsxvier $  &  $ \aeinsxfuenf $  &  $ \aeinsxsechs $ \\

\hline  \leitspaltezwei  &  $ \azweixeins $  &  $ \azweixzwei $  &  $ \azweixdrei $  &  $ \azweixvier $  &  $ \azweixfuenf $  &  $ \azweixsechs $   \\

\hline  \leitspaltedrei  &  $ \adreixeins $  &  $ \adreixzwei $  &  $ \adreixdrei $  &  $ \adreixvier $  &  $ \adreixfuenf $   &  $ \adreixsechs $   \\

\hline  \leitspaltevier  &  $ \avierxeins $  &  $ \avierxzwei $  &  $ \avierxdrei $  &  $ \avierxvier $  &  $ \avierxfuenf $   &  $ \avierxsechs $  \\

\hline  \leitspaltefuenf  &  $ \afuenfxeins $  &  $ \afuenfxzwei $  &  $ \afuenfxdrei $  &  $ \afuenfxvier $  &  $ \afuenfxfuenf $ &  $ \afuenfxsechs $   \\

\hline  \leitspaltesechs  &  $ \asechsxeins $  &  $ \asechsxzwei $  &  $ \asechsxdrei $  &  $ \asechsxvier $  &  $ \asechsxfuenf $  &  $ \asechsxsechs $  \\

\hline

\end{tabular}

\end{center} }

%Tabellen mit sieben Zeilen

\newcommand{\tabelleleitsiebenxsieben}{

\begin{center}

\begin{tabular}{|c||c|c|c|c|c|c|c|}

\hline  \leitzeilenull  &  \leitzeileeins  &  \leitzeilezwei  &  \leitzeiledrei  &  \leitzeilevier  &  \leitzeilefuenf  &  \leitzeilesechs &  \leitzeilesieben  \\

\hline \hline  \leitspalteeins  &  $ \aeinsxeins $  &  $ \aeinsxzwei $  &  $ \aeinsxdrei $  &  $ \aeinsxvier $  &  $ \aeinsxfuenf $  &  $ \aeinsxsechs $ &  $ \aeinsxsieben $ \\

\hline  \leitspaltezwei  &  $ \azweixeins $  &  $ \azweixzwei $  &  $ \azweixdrei $  &  $ \azweixvier $  &  $ \azweixfuenf $  &  $ \azweixsechs $  &  $ \azweixsieben $ \\

\hline  \leitspaltedrei  &  $ \adreixeins $  &  $ \adreixzwei $  &  $ \adreixdrei $  &  $ \adreixvier $  &  $ \adreixfuenf $   &  $ \adreixsechs $ &  $ \adreixsieben $  \\

\hline  \leitspaltevier  &  $ \avierxeins $  &  $ \avierxzwei $  &  $ \avierxdrei $  &  $ \avierxvier $  &  $ \avierxfuenf $   &  $ \avierxsechs $ &  $ \avierxsieben $ \\

\hline  \leitspaltefuenf  &  $ \afuenfxeins $  &  $ \afuenfxzwei $  &  $ \afuenfxdrei $  &  $ \afuenfxvier $  &  $ \afuenfxfuenf $ &  $ \afuenfxsechs $ &  $ \afuenfxsieben $  \\

\hline  \leitspaltesechs  &  $ \asechsxeins $  &  $ \asechsxzwei $  &  $ \asechsxdrei $  &  $ \asechsxvier $  &  $ \asechsxfuenf $  &  $ \asechsxsechs $ &  $ \asechsxsieben $ \\

\hline  \leitspaltesieben  &  $ \asiebenxeins $  &  $ \asiebenxzwei $  &  $ \asiebenxdrei $  &  $ \asiebenxvier $  &  $ \asiebenxfuenf $  &  $ \asiebenxsechs $ &  $ \asiebenxsieben $ \\

\hline

\end{tabular}

\end{center} }

%Tabellen mit acht Zeilen

\newcommand{\tabelleleitachtxdrei}{

\begin{center}

\begin{tabular}{|c||c|c|c|}

\hline  \leitzeilenull  &  \leitzeileeins  &  \leitzeilezwei  &  \leitzeiledrei     \\

\hline \hline  \leitspalteeins  &  $ \aeinsxeins $  &  $ \aeinsxzwei $  &  $ \aeinsxdrei $    \\

\hline  \leitspaltezwei  &  $ \azweixeins $  &  $ \azweixzwei $  &  $ \azweixdrei $     \\

\hline  \leitspaltedrei  &  $ \adreixeins $  &  $ \adreixzwei $  &  $ \adreixdrei $    \\

\hline  \leitspaltevier  &  $ \avierxeins $  &  $ \avierxzwei $  &  $ \avierxdrei $      \\

\hline  \leitspaltefuenf  &  $ \afuenfxeins $  &  $ \afuenfxzwei $  &  $ \afuenfxdrei $     \\

\hline  \leitspaltesechs  &  $ \asechsxeins $  &  $ \asechsxzwei $  &  $ \asechsxdrei $      \\

\hline  \leitspaltesieben  &  $ \asiebenxeins $  &  $ \asiebenxzwei $  &  $ \asiebenxdrei $     \\

\hline  \leitspalteacht  &  $ \aachtxeins $  &  $ \aachtxzwei $  &  $ \aachtxdrei $      \\

\hline

\end{tabular}

\end{center} }

%Tabellen mit zehn Zeilen

\newcommand{\tabelleleitzehnxdrei}{

\begin{center}

\begin{tabular}{|c||c|c|c|}

\hline  \leitzeilenull  &  \leitzeileeins  &  \leitzeilezwei  &  \leitzeiledrei     \\

\hline \hline  \leitspalteeins  &  $ \aeinsxeins $  &  $ \aeinsxzwei $  &  $ \aeinsxdrei $    \\

\hline  \leitspaltezwei  &  $ \azweixeins $  &  $ \azweixzwei $  &  $ \azweixdrei $     \\

\hline  \leitspaltedrei  &  $ \adreixeins $  &  $ \adreixzwei $  &  $ \adreixdrei $    \\

\hline  \leitspaltevier  &  $ \avierxeins $  &  $ \avierxzwei $  &  $ \avierxdrei $      \\

\hline  \leitspaltefuenf  &  $ \afuenfxeins $  &  $ \afuenfxzwei $  &  $ \afuenfxdrei $     \\

\hline  \leitspaltesechs  &  $ \asechsxeins $  &  $ \asechsxzwei $  &  $ \asechsxdrei $      \\

\hline  \leitspaltesieben  &  $ \asiebenxeins $  &  $ \asiebenxzwei $  &  $ \asiebenxdrei $     \\

\hline  \leitspalteacht  &  $ \aachtxeins $  &  $ \aachtxzwei $  &  $ \aachtxdrei $      \\

\hline  \leitspalteneun  &  $ \aneunxeins $  &  $ \aneunxzwei $  &  $ \aneunxdrei $      \\

\hline  \leitspaltezehn  &  $ \azehnxeins $  &  $ \azehnxzwei $  &  $ \azehnxdrei $      \\

\hline

\end{tabular}

\end{center} }

%Tabellen mit elf Zeilen

\newcommand{\tabelleleitelfxvier}{

\begin{center}

\begin{tabular}{|c||c|c|c|c|}

\hline  \leitzeilenull  &  \leitzeileeins  &  \leitzeilezwei  &  \leitzeiledrei  &  \leitzeilevier  \\

\hline \hline  \leitspalteeins  &  $ \aeinsxeins $  &  $ \aeinsxzwei $  &  $ \aeinsxdrei $  &  $ \aeinsxvier $  \\

\hline  \leitspaltezwei  &  $ \azweixeins $  &  $ \azweixzwei $  &  $ \azweixdrei $  &  $ \azweixvier $   \\

\hline  \leitspaltedrei  &  $ \adreixeins $  &  $ \adreixzwei $  &  $ \adreixdrei $  &  $ \adreixvier $   \\

\hline  \leitspaltevier  &  $ \avierxeins $  &  $ \avierxzwei $  &  $ \avierxdrei $  &  $ \avierxvier $   \\

\hline  \leitspaltefuenf  &  $ \afuenfxeins $  &  $ \afuenfxzwei $  &  $ \afuenfxdrei $  &  $ \afuenfxvier $   \\

\hline  \leitspaltesechs  &  $ \asechsxeins $  &  $ \asechsxzwei $  &  $ \asechsxdrei $  &  $ \asechsxvier $   \\

\hline  \leitspaltesieben  &  $ \asiebenxeins $  &  $ \asiebenxzwei $  &  $ \asiebenxdrei $  &  $ \asiebenxvier $    \\

\hline  \leitspalteacht  &  $ \aachtxeins $  &  $ \aachtxzwei $  &  $ \aachtxdrei $  &  $ \aachtxvier $   \\

\hline  \leitspalteneun  &  $ \aneunxeins $  &  $ \aneunxzwei $  &  $ \aneunxdrei $  &  $ \aneunxvier $   \\

\hline  \leitspaltezehn  &  $ \azehnxeins $  &  $ \azehnxzwei $  &  $ \azehnxdrei $  &  $ \azehnxvier $   \\

\hline  \leitspalteelf  &  $ \aelfxeins $  &  $ \aelfxzwei $  &  $ \aelfxdrei $  &  $ \aelfxvier $   \\

\hline

\end{tabular}

\end{center} }

%Tabellen mit zwölf Zeilen

\newcommand{\tabelleleitzwoelfxsieben}{

\begin{center}

\begin{tabular}{|c||c|c|c|c|c|c|}

\hline  \leitzeilenull  &  \leitzeileeins  &  \leitzeilezwei  &  \leitzeiledrei  &  \leitzeilevier &  \leitzeilefuenf  &  \leitzeilesechs &  \leitzeilesieben \\

\hline \hline  \leitspalteeins  &  $ \aeinsxeins $  &  $ \aeinsxzwei $  &  $ \aeinsxdrei $  &  $ \aeinsxvier $  &  $ \aeinsxfuenf $  &  $ \aeinsxsechs $  &  $ \aeinsxsieben $  \\

\hline  \leitspaltezwei  &  $ \azweixeins $  &  $ \azweixzwei $  &  $ \azweixdrei $  &  $ \azweixvier $  &  $ \azweixfuenf $  &  $ \azweixsechs $  &  $ \azweixsieben $   \\

\hline  \leitspaltedrei  &  $ \adreixeins $  &  $ \adreixzwei $  &  $ \adreixdrei $  &  $ \adreixvier $  &  $ \adreixfuenf $  &  $ \adreixsechs $  &  $ \adreixsieben $   \\

\hline  \leitspaltevier  &  $ \avierxeins $  &  $ \avierxzwei $  &  $ \avierxdrei $  &  $ \avierxvier $   &  $ \avierxfuenf $  &  $ \avierxsechs $  &  $ \avierxsieben $  \\

\hline  \leitspaltefuenf  &  $ \afuenfxeins $  &  $ \afuenfxzwei $  &  $ \afuenfxdrei $  &  $ \afuenfxvier $ &  $ \afuenfxfuenf $  &  $ \afuenfxsechs $  &  $ \afuenfxsieben $  \\

\hline  \leitspaltesechs  &  $ \asechsxeins $  &  $ \asechsxzwei $  &  $ \asechsxdrei $  &  $ \asechsxvier $  &  $ \asechsxfuenf $  &  $ \asechsxsechs $  &  $ \asechsxsieben $  \\

\hline  \leitspaltesieben  &  $ \asiebenxeins $  &  $ \asiebenxzwei $  &  $ \asiebenxdrei $  &  $ \asiebenxvier $  &  $ \asiebenxfuenf $  &  $ \asiebenxsechs $  &  $ \asiebenxsieben $   \\

\hline  \leitspalteacht  &  $ \aachtxeins $  &  $ \aachtxzwei $  &  $ \aachtxdrei $  &  $ \aachtxvier $  &  $ \aachtxfuenf $  &  $ \aachtxsechs $  &  $ \aachtxsieben $ \\

\hline  \leitspalteneun  &  $ \aneunxeins $  &  $ \aneunxzwei $  &  $ \aneunxdrei $  &  $ \aneunxvier $  &  $ \aneunxfuenf $  &  $ \aneunxsechs $  &  $ \aneunxsieben $  \\

\hline  \leitspaltezehn  &  $ \azehnxeins $  &  $ \azehnxzwei $  &  $ \azehnxdrei $  &  $ \azehnxvier $ &  $ \azehnxfuenf $  &  $ \azehnxsechs $  &  $ \azehnxsieben $   \\

\hline  \leitspalteelf  &  $ \aelfxeins $  &  $ \aelfxzwei $  &  $ \aelfxdrei $  &  $ \aelfxvier $  &  $ \aelfxfuenf $  &  $ \aelfxsechs $  &  $ \aelfxsieben $  \\

\hline  \leitspaltezwoelf  &  $ \azwoelfxeins $  &  $ \azwoelfxzwei $  &  $ \azwoelfxdrei $  &  $ \zwoelfxvier $  &  $ \azwoelfxfuenf $  &  $ \azwoelfxsechs $  &  $ \zwoelfxsieben $ \\

\hline

\end{tabular}

\end{center} }

%Tabellen mit sechzehn Zeilen

\newcommand{\tabelleleitsechzehnxzwei}{

\begin{center}

\begin{tabular}{|c||c|c|}

\hline  \leitzeilenull  &  \leitzeileeins  &  \leitzeilezwei   \\

\hline \hline  \leitspalteeins  &  $ \aeinsxeins $  &  $ \aeinsxzwei $ \\

\hline  \leitspaltezwei  &  $ \azweixeins $  &  $ \azweixzwei $  \\

\hline  \leitspaltedrei  &  $ \adreixeins $  &  $ \adreixzwei $  \\

\hline  \leitspaltevier  &  $ \avierxeins $  &  $ \avierxzwei $  \\

\hline  \leitspaltefuenf  &  $ \afuenfxeins $  &  $ \afuenfxzwei $  \\

\hline  \leitspaltesechs  &  $ \asechsxeins $  &  $ \asechsxzwei $  \\

\hline  \leitspaltesieben  &  $ \asiebenxeins $  &  $ \asiebenxzwei $  \\

\hline  \leitspalteacht  &  $ \aachtxeins $  &  $ \aachtxzwei $  \\

\hline  \leitspalteneun  &  $ \aneunxeins $  &  $ \aneunxzwei $   \\

\hline  \leitspaltezehn  &  $ \azehnxeins $  &  $ \azehnxzwei $  \\

\hline  \leitspalteelf  &  $ \aelfxeins $  &  $ \aelfxzwei $   \\

\hline  \leitspaltezwoelf  &  $ \azwoelfxeins $  &  $ \azwoelfxzwei $   \\

\hline  \leitspaltedreizehn  &  $ \adreizehnxeins $  &  $ \adreizehnxzwei $   \\

\hline  \leitspaltevierzehn  &  $ \avierzehnxeins $  &  $ \avierzehnxzwei $  \\

\hline  \leitspaltefuenfzehn  &  $ \afuenfzehnxeins $  &  $ \afuenfzehnxzwei $   \\

\hline  \leitspaltesechzehn  &  $ \asechzehnxeins $  &  $ \asechzehnxzwei $   \\

\hline

\end{tabular}

\end{center} }

\newcommand{\tabelleleitsechzehnxsechs}{

\begin{center}

\begin{tabular}{|c||c|c|c|c|c|c|}

\hline  \leitzeilenull  &  \leitzeileeins  &  \leitzeilezwei  &  \leitzeiledrei  &  \leitzeilevier  &  \leitzeilefuenf  &  \leitzeilesechs  \\

\hline \hline  \leitspalteeins  &  $ \aeinsxeins $  &  $ \aeinsxzwei $  &  $ \aeinsxdrei $  &  $ \aeinsxvier $  &  $ \aeinsxfuenf $  &  $ \aeinsxsechs $  \\

\hline  \leitspaltezwei  &  $ \azweixeins $  &  $ \azweixzwei $  &  $ \azweixdrei $  &  $ \azweixvier $  &  $ \azweixfuenf $  &  $ \azweixsechs $  \\

\hline  \leitspaltedrei  &  $ \adreixeins $  &  $ \adreixzwei $  &  $ \adreixdrei $  &  $ \adreixvier $  &  $ \adreixfuenf $  &  $ \adreixsechs $  \\

\hline  \leitspaltevier  &  $ \avierxeins $  &  $ \avierxzwei $  &  $ \avierxdrei $  &  $ \avierxvier $  &  $ \avierxfuenf $  &  $ \avierxsechs $  \\

\hline  \leitspaltefuenf  &  $ \afuenfxeins $  &  $ \afuenfxzwei $  &  $ \afuenfxdrei $  &  $ \afuenfxvier $  &  $ \afuenfxfuenf $  &  $ \afuenfxsechs $  \\

\hline  \leitspaltesechs  &  $ \asechsxeins $  &  $ \asechsxzwei $  &  $ \asechsxdrei $  &  $ \asechsxvier $  &  $ \asechsxfuenf $  &  $ \asechsxsechs $  \\

\hline  \leitspaltesieben  &  $ \asiebenxeins $  &  $ \asiebenxzwei $  &  $ \asiebenxdrei $  &  $ \asiebenxvier $  &  $ \asiebenxfuenf $  &  $ \asiebenxsechs $  \\

\hline  \leitspalteacht  &  $ \aachtxeins $  &  $ \aachtxzwei $  &  $ \aachtxdrei $  &  $ \aachtxvier $  &  $ \aachtxfuenf $  &  $ \aachtxsechs $  \\

\hline  \leitspalteneun  &  $ \aneunxeins $  &  $ \aneunxzwei $  &  $ \aneunxdrei $  &  $ \aneunxvier $  &  $ \aneunxfuenf $  &  $ \aneunxsechs $  \\

\hline  \leitspaltezehn  &  $ \azehnxeins $  &  $ \azehnxzwei $  &  $ \azehnxdrei $  &  $ \azehnxvier $  &  $ \azehnxfuenf $  &  $ \azehnxsechs $  \\

\hline  \leitspalteelf  &  $ \aelfxeins $  &  $ \aelfxzwei $  &  $ \aelfxdrei $  &  $ \aelfxvier $  &  $ \aelfxfuenf $  &  $ \aelfxsechs $  \\

\hline  \leitspaltezwoelf  &  $ \azwoelfxeins $  &  $ \azwoelfxzwei $  &  $ \azwoelfxdrei $  &  $ \azwoelfxvier $  &  $ \azwoelfxfuenf $  &  $ \azwoelfxsechs $  \\

\hline  \leitspaltedreizehn  &  $ \adreizehnxeins $  &  $ \adreizehnxzwei $  &  $ \adreizehnxdrei $  &  $ \adreizehnxvier $  &  $ \adreizehnxfuenf $  &  $ \adreizehnxsechs $  \\

\hline  \leitspaltevierzehn  &  $ \avierzehnxeins $  &  $ \avierzehnxzwei $  &  $ \avierzehnxdrei $  &  $ \avierzehnxvier $  &  $ \avierzehnxfuenf $  &  $ \avierzehnxsechs $  \\

\hline  \leitspaltefuenfzehn  &  $ \afuenfzehnxeins $  &  $ \afuenfzehnxzwei $  &  $ \afuenfzehnxdrei $  &  $ \afuenfzehnxvier $  &  $ \afuenfzehnxfuenf $  &  $ \afuenfzehnxsechs $  \\

\hline  \leitspaltesechzehn  &  $ \asechzehnxeins $  &  $ \asechzehnxzwei $  &  $ \asechzehnxdrei $  &  $ \asechzehnxvier $  &  $ \asechzehnxfuenf $  &  $ \asechzehnxsechs $  \\

\hline

\end{tabular}

\end{center} }

%Tabellen mit achtzehn Zeilen

\newcommand{\tabelleleitachtzehnxneun}{

\begin{center}

\begin{tabular}{|c||c|c|c|c|c|c|c|c|c|}

\hline  \leitzeilenull  &  \leitzeileeins  &  \leitzeilezwei  &  \leitzeiledrei  &  \leitzeilevier  &  \leitzeilefuenf  &  \leitzeilesechs &  \leitzeilesieben  &  \leitzeileacht  &  \leitzeileneun \\

\hline \hline  \leitspalteeins  &  $ \aeinsxeins $  &  $ \aeinsxzwei $  &  $ \aeinsxdrei $  &  $ \aeinsxvier $  &  $ \aeinsxfuenf $  &  $ \aeinsxsechs $ &  $ \aeinsxsieben $  &  $ \aeinsxacht $  &  $ \aeinsxneun $  \\

\hline  \leitspaltezwei  &  $ \azweixeins $  &  $ \azweixzwei $  &  $ \azweixdrei $  &  $ \azweixvier $  &  $ \azweixfuenf $  &  $ \azweixsechs $ &  $ \azweixsieben $  &  $ \azweixacht $  &  $ \azweixneun $  \\

\hline  \leitspaltedrei  &  $ \adreixeins $  &  $ \adreixzwei $  &  $ \adreixdrei $  &  $ \adreixvier $  &  $ \adreixfuenf $  &  $ \adreixsechs $ &  $ \adreixsieben $  &  $ \adreixacht $  &  $ \adreixneun $  \\

\hline  \leitspaltevier  &  $ \avierxeins $  &  $ \avierxzwei $  &  $ \avierxdrei $  &  $ \avierxvier $  &  $ \avierxfuenf $  &  $ \avierxsechs $  &  $ \avierxsieben $  &  $ \avierxacht $  &  $ \avierxneun $ \\

\hline  \leitspaltefuenf  &  $ \afuenfxeins $  &  $ \afuenfxzwei $  &  $ \afuenfxdrei $  &  $ \afuenfxvier $  &  $ \afuenfxfuenf $  &  $ \afuenfxsechs $ &  $ \afuenfxsieben $  &  $ \afuenfxacht $  &  $ \afuenfxneun $  \\

\hline  \leitspaltesechs  &  $ \asechsxeins $  &  $ \asechsxzwei $  &  $ \asechsxdrei $  &  $ \asechsxvier $  &  $ \asechsxfuenf $  &  $ \asechsxsechs $ &  $ \asechsxsieben $  &  $ \asechsxacht $  &  $ \asechsxneun $  \\

\hline  \leitspaltesieben  &  $ \asiebenxeins $  &  $ \asiebenxzwei $  &  $ \asiebenxdrei $  &  $ \asiebenxvier $  &  $ \asiebenxfuenf $  &  $ \asiebenxsechs $ &  $ \asiebenxsieben $  &  $ \asiebenxacht $  &  $ \asiebenxneun $  \\

\hline  \leitspalteacht  &  $ \aachtxeins $  &  $ \aachtxzwei $  &  $ \aachtxdrei $  &  $ \aachtxvier $  &  $ \aachtxfuenf $  &  $ \aachtxsechs $ &  $ \aachtxsieben $  &  $ \aachtxacht $  &  $ \aachtxneun $  \\

\hline  \leitspalteneun  &  $ \aneunxeins $  &  $ \aneunxzwei $  &  $ \aneunxdrei $  &  $ \aneunxvier $  &  $ \aneunxfuenf $  &  $ \aneunxsechs $ &  $ \aneunxsieben $  &  $ \aneunxacht $  &  $ \aneunxneun $  \\

\hline  \leitspaltezehn  &  $ \azehnxeins $  &  $ \azehnxzwei $  &  $ \azehnxdrei $  &  $ \azehnxvier $  &  $ \azehnxfuenf $  &  $ \azehnxsechs $ &  $ \azehnxsieben $  &  $ \azehnxacht $  &  $ \azehnxneun $  \\

\hline  \leitspalteelf  &  $ \aelfxeins $  &  $ \aelfxzwei $  &  $ \aelfxdrei $  &  $ \aelfxvier $  &  $ \aelfxfuenf $  &  $ \aelfxsechs $ &  $ \aelfxsieben $  &  $ \aelfxacht $  &  $ \aelfxneun $  \\

\hline  \leitspaltezwoelf  &  $ \azwoelfxeins $  &  $ \azwoelfxzwei $  &  $ \azwoelfxdrei $  &  $ \azwoelfxvier $  &  $ \azwoelfxfuenf $  &  $ \azwoelfxsechs $ &  $ \azwoelfxsieben $  &  $ \azwoelfxacht $  &  $ \azwoelfxneun $  \\

\hline  \leitspaltedreizehn  &  $ \adreizehnxeins $  &  $ \adreizehnxzwei $  &  $ \adreizehnxdrei $  &  $ \adreizehnxvier $  &  $ \adreizehnxfuenf $  &  $ \adreizehnxsechs $ &  $ \adreizehnxsieben $  &  $ \adreizehnxacht $  &  $ \adreizehnxneun $  \\

\hline  \leitspaltevierzehn  &  $ \avierzehnxeins $  &  $ \avierzehnxzwei $  &  $ \avierzehnxdrei $  &  $ \avierzehnxvier $  &  $ \avierzehnxfuenf $  &  $ \avierzehnxsechs $  &  $ \avierzehnxsieben $  &  $ \avierzehnxacht $  &  $ \avierzehnxneun $ \\

\hline  \leitspaltefuenfzehn  &  $ \afuenfzehnxeins $  &  $ \afuenfzehnxzwei $  &  $ \afuenfzehnxdrei $  &  $ \afuenfzehnxvier $  &  $ \afuenfzehnxfuenf $  &  $ \afuenfzehnxsechs $ &  $ \afuenfzehnxsieben $  &  $ \afuenfzehnxacht $  &  $ \afuenfzehnxneun $  \\

\hline  \leitspaltesechzehn  &  $ \asechzehnxeins $  &  $ \asechzehnxzwei $  &  $ \asechzehnxdrei $  &  $ \asechzehnxvier $  &  $ \asechzehnxfuenf $  &  $ \asechzehnxsechs $ &  $ \asechzehnxsieben $  &  $ \asechzehnxacht $  &  $ \asechzehnxneun $  \\

\hline  \leitspaltesiebzehn  &  $ \asiebzehnxeins $  &  $ \asiebzehnxzwei $  &  $ \asiebzehnxdrei $  &  $ \asiebzehnxvier $  &  $ \asiebzehnxfuenf $  &  $ \asiebzehnxsechs $ &  $ \asiebzehnxsieben $  &  $ \asiebzehnxacht $  &  $ \asiebzehnxneun $  \\

\hline  \leitspalteachtzehn  &  $ \aachtzehnxeins $  &  $ \aachtzehnxzwei $  &  $ \aachtzehnxdrei $  &  $ \aachtzehnxvier $  &  $ \aachtzehnxfuenf $  &  $ \aachtzehnxsechs $ &  $ \aachtzehnxsieben $  &  $ \aachtzehnxacht $  &  $ \aachtzehnxneun $  \\

\hline

\end{tabular}

\end{center} }

\newcommand{\bruchhilfealign}{\cr & & }

\newcommand{\bruchhilfedisp}{\]  \[ }

\newcommand{\bruchhilfe}{ \- }

%[[Projekt:Semantische Vorlagen/Notlösungen/Latex]]

%
%Bei verschachtelten ifthen-Befehlen
%kann es Probleme geben, daher
%

\newcommand{\faktuebergangpos}[1]{ \faktuebergangskip  {\hspace{-0,55cm} } {#1} }

%Ende des Grundvorspannes

%Globale Strukturierungen

\newcommand{\seitenueberschrift}[1]{\seitenueberschriftvorskip
\begin{center} \large{ #1} \end{center} \seitenueberschriftnachskip}

\newcommand{\zwischenueberschrift}[1]{\section*{#1} }

\newcommand{\zwischenueberschriftaufgaben}[1]{ \zwischenueberschriftvorskip \begin{center} {\bf #1} \end{center} \zwischenueberschriftnachskip}

\newcommand{\unterueberschrift}[1]{ \unterueberschriftvorskip \hspace{0.5cm} {\em #1}  \unterueberschriftnachskip}

\renewcommand{\definitionbenennung}[1]{}

%Grund:Bei Definitionen im Normaltext
%nicht nötig, da Begriff früh kommt

\renewcommand{\beispielbenennung}[1]{}

\renewcommand{\bemerkungbenennung}[1]{}

%Umbrüche bei großer Schrift (entfällt)

%[[Projekt:Semantische Vorlagen/Mathematischer Modus/Umbrüche/Latex]]

\newcommand{\mathdisplaybruch}{}

\newcommand{\mathdisplaybruchinklammer}{}

\newcommand{\mathl}[2]{$#1$#2}

\newcommand{\mathlk}[2]{$#1$#2}

%Nicht durch semantische Vorlage produzierte Befehle

\newcommand{\mathdisplaybruchbruch}{ \]  \[ }

%Gestaltung der vertikalen und horizontalen Abstände

%[[Projekt:Semantische Vorlagen/Latex/Vorlesungsskript/Skips]]

%Abstände bei Überschriften

\newcommand{\seitenueberschriftvorskip}{\par \bigskip \bigskip \bigskip }

\newcommand{\seitenueberschriftnachskip}{\par \bigskip}

\newcommand{\zwischenueberschriftvorskip}{\par \bigskip}

\newcommand{\zwischenueberschriftnachskip}{\par \smallskip}

\newcommand{\unterueberschriftvorskip}{\par \bigskip}

\newcommand{\unterueberschriftnachskip}{\par \smallskip}

%Aufzählungen und Listen

\newcommand{\listskip}{\smallskip}

%Abstände bei mathematischen Strukturen

\newcommand{\aufgabevorskip}{\bigskip}

\newcommand{\aufgabepunktskip}{\par}

\newcommand{\aufgabenachskip}{}

\newcommand{\aufgabeloesungskip}{\par \bigskip}

\newcommand{\beispielvorskip}{}

\newcommand{\beispielnachskip}{}

\newcommand{\beispielbenennungnachskip}{}

\newcommand{\bemerkungvorskip}{}

\newcommand{\bemerkungnachskip}{}

\newcommand{\bemerkungbenennungnachskip}{}

\newcommand{\definitionvorskip}{}

\newcommand{\definitionnachskip}{}

\newcommand{\definitionbenennungnachskip}{}

\newcommand{\faktvorskip}{}

\newcommand{\faktnachskip}{}

\newcommand{\faktbenennungvorskip}{}

\newcommand{\faktbenennungnachskip}{}

\newcommand{\faktsituationskip}{}

\newcommand{\faktvoraussetzungskip}{}

\newcommand{\faktuebergangskip}{}

\newcommand{\faktfolgerungskip}{}

\newcommand{\faktzusatzskip}{}

%Sonstiges

\newcommand{\deckblattabstand}{\vspace{2cm}}

\newcommand{\gesamtskriptnewpage}{ }

\newcommand{\einzeltextnewpage}{ }

%[[Projekt:Semantische Vorlagen/Latex/Vorlesungsskript/Horizontale Abstände]]

\newcommand{\leerzeichen}{ }

%Seitengestaltung

%[[Projekt:Semantische Vorlagen/Latex/Vorlesungsskript/Seitengestaltung]]

%Einzelne Abstände

% page geometry supplied by reader wrapper
% page geometry supplied by reader wrapper

% page geometry supplied by reader wrapper
% page geometry supplied by reader wrapper

% page geometry supplied by reader wrapper
% page geometry supplied by reader wrapper

\setlength{\parindent}{0cm}

\setlength{\parskip}{1ex}

%Bildgestaltung

%[[Projekt:Semantische Vorlagen/Latex/Vorlesungsskript/Bildgestaltung]]

\newcommand{\bild}[1]{\vspace{4mm} #1}

\newcommand{\bildtext}[1]{\begin{center} {\small #1 } \end{center} }

%Lokale Einstellungen

%\input{Lokalername}

\newcommand{\bildeinlesung}[1]{../authority/media/#1}

\renewcommand{\bildeinlesung}[2]{../authority/media/#1.#2}

\newcommand{\bildeinlesungpng}[2]{../authority/media/#1.png}

\newcommand{\bildeinlesungPNG}[2]{../authority/media/#1.png}

\newcommand{\bildeinlesungjpg}[2]{../authority/media/#1.jpg}

\newcommand{\bildeinlesungJPG}[2]{../authority/media/#1.jpg}

\newcommand{\bildeinlesungjpeg}[2]{../authority/media/#1.jpg}

\newcommand{\bildeinlesungsvg}[2]{generated/media/#1.png}

\newcommand{\bildeinlesunggif}[2]{../authority/media/#1.png}

\newcommand{\bildeinlesungGIF}[2]{../authority/media/#1.png}

\newcommand{\bildeinlesungxcf}[2]{../authority/media/#1.png}

%Klausurgestaltung

%[[Projekt:Semantische Vorlagen/Klausurgestaltung/Latex]]

\newcommand{\fachbereich}{Fachbereich}

\newcommand{\dozent}{Dozent}

\newcommand{\klausurdatum}{Datum}

\newcommand{\klausurgebiet}{Mathematik}

\newcommand{\klausurtyp}{Klausur}

\newcommand{\klausurmodus}{Dauer: Zehn Minuten Einlesezeit und zwei volle Stunden Bearbeitungszeit. Es sind keine Hilfsmittel erlaubt. Alle Antworten sind zu begr\"unden.

Es gibt insgesamt 64 Punkte. Es gilt die Sockelregelung, d.h. die Bewertung pro Aufgabe(nteil) beginnt in der Regel bei der halben Punktzahl. Zum Bestehen braucht man $16$ Punkte, ab $32$ Punkten gibt es eine Eins.

Tragen Sie auf dem Deckblatt und jedem weiteren Blatt Ihren Namen und Ihre
Matrikelnummer leserlich ein. Viel Erfolg!}

\newcommand{\klausurname}{
\renewcommand{\arraystretch}{1.5}
\begin{tabular}{@{}p{3.2cm}p{10cm}}
Name, Vorname: &
..................................................................................\\
Matrikelnummer: &
..................................................................................\\
\end{tabular}
}

\newcommand{\klausurvorspann}[5]{\textbf{\fachbereich}\hfill \textbf{\klausurdatum}
\\
\textbf{\dozent} \hfill
\vspace{4ex}

\begin{center}
{\bfseries{\large \klausurgebiet \vspace{2ex}

\klausurtyp} }
\end{center}

\vspace{2ex}

\klausurmodus

\vspace{2ex}

\klausurname
}

\newcommand{\klausurnote}{\begin{tabular}{@{}p{3.2cm}p{10cm}} Note: & \end{tabular} }

\newcommand{\klausurtaeuschungwarnung}{ \vspace{2ex}{Mir ist bewu\ss t, dass bei einem schweren T\"auschungsversuch s\"amtliche Pr\"u\-fungsvorleis\-tungen in diesem Kurs verfallen. Ein schwerer T\"auschungsversuch liegt insbesondere beim Einsatz von Literatur und von elektronischen Hilfsmitteln vor. \vspace{1ex} } }

\newcommand{\klausureinverstaendnis}{
\vspace{2ex}
Ich erkl\"are mich durch meine Unterschrift einverstanden, dass mein Klausurergebnis mit meiner Matrikelnummer im Internet bekanntgegeben wird.
\vspace{1ex}

\begin{tabular}{@{}p{3.2cm}p{10cm}}
Unterschrift: &
.....................................................................................
\end{tabular}
}

%Variante für Testklausur

\newcommand{\testklausurvorspann}[5]{\textbf{\fachbereich}\hfill \textbf{\klausurdatum} \\ \textbf{\dozent} \hfill \vspace{4ex}

\begin{center} {\bfseries{\large \klausurgebiet \vspace{2ex}

\klausurtyp} } \end{center}

\vspace{2ex}

\testklausurmodus

\vspace{3ex}

\klausurname }

\newcommand{\testklausurmodus}{Dauer: Zehn Minuten Einlesezeit und zwei volle Stunden Bearbeitungszeit.

Es sind keine Hilfsmittel erlaubt.

Alle Antworten sind zu begr\"unden.

Es gibt insgesamt 64 Punkte. Es gilt die Sockelregelung, d.h. die Bewertung pro Aufgabe(nteil) beginnt in der Regel bei der halben Punktzahl.

Zur Orientierung: Zum Bestehen braucht man $16$ Punkte, ab $32$ Punkten gibt es eine Eins.

Tragen Sie auf dem Deckblatt und jedem weiteren Blatt Ihren Namen und Ihre Matrikelnummer leserlich ein.

Viel Erfolg!}

%\renewcommand{\inputaufgabegibtloesung}[4]{\aufgabevorskip \begin{aufgabe} \punkte {#1} \aufgabepunktskip #2 #3 #4\end{aufgabe} \aufgabenachskip}

\newcommand{\klausurmitloesungenvorspann}[5]{\textbf{\fachbereich}\hfill \textbf{\klausurdatum} \\ \textbf{\dozent} \hfill \vspace{4ex}

\begin{center} {\bfseries{\large \klausurgebiet \vspace{2ex}

\klausurtyp \hspace{0mm} mit L\"osungen} } \end{center}

}

%Lizenzierung und Bildverzeichnis

%[[Projekt:Semantische Vorlagen/Latex/Vorlesungsskript/Lizenzierung]]

\newcommand{\Lizenztext}{Teks sumber dibuat oleh Holger Brenner di Wikiversity berbahasa Jerman dan digunakan di sini berdasarkan CC BY-SA 4.0. Media mengikuti lisensi per berkas.}

\newcommand{\Abbildungslizenzierung}{ Erl\"auterung: Die in diesem Text verwendeten Bilder stammen aus Commons (also von http://commons.wikimedia.org) und haben eine Lizenz, die die Verwendung hier erlaubt. Die Bilder werden mit ihren Dateinamen auf Commons angef\"uhrt zusammen mit ihrem Autor bzw. Hochlader und der Lizenz. }

%Konfiguration des Bildverzeichnisses

\newcommand{\bildlizenz}[6]{
\ifthenelse {\equal{#2}{}}
{\addcontentsline{lof}{figure}{Sumber = #1, Kreator = Pengguna #3 di #4, Lisensi = #5 \bildlizenzskip  }}
{ \ifthenelse {\equal{#3}{}}
{ \addcontentsline{lof}{figure}{ Sumber = #1, Kreator = #2, Lisensi = #5 \bildlizenzskip }}
{ \addcontentsline{lof}{figure}{ Sumber = #1, Kreator = #2 (diunggah oleh Pengguna #3 di #4), Lisensi = #5 \bildlizenzskip }} } }

\newcommand{\bildlizenzskip}{ }

%Englische Variante

\newcommand{\ignoriereeng}[1]{}

\ignoriereeng{

%[[Projekt:Semantische Vorlagen/Latex/Englische Zusätze]]

\newtheorem{corollary}[fakt]{Corollary}

\newtheorem{Corollary}[fakt]{Corollary}

\theoremstyle{definition}

\newtheorem{example}[fakt]{Example}

\newtheorem{remark}[fakt]{Remark}

\newcommand{\inputremark}[2] {\bemerkungvorskip \begin{remark} \bemerkungbenennung{#1} \bemerkungbenennungnachskip #2 \end{remark} \bemerkungnachskip}

\newcommand{\inputexample}[2] {\beispielvorskip \begin{example} \beispielbenennung{#1} \beispielbenennungnachskip #2 \end{example} \beispielnachskip}

\newcommand{\inputfaktproof}[5] {\faktvorskip \begin{#2}

%\label{#1} (Absetzen, sonst kann das folgende dahinter rutschen)

\faktbenennung{#3} \faktbenennungnachskip #4 \end{#2} \faktnachskip \beweis{#5}}

\renewcommand{\proofname}{\hspace{-0.64cm}{\it Proof}}

}

%Variante für Artikel

\newcommand{\ignoriereart}[1]{}

\ignoriereart{
\newcommand{\titelueberschrift}[1]{\title[#1] { #1} \maketitle }

\renewcommand{\seitenueberschrift}[1]{ \section {#1} }

\renewcommand{\zwischenueberschrift}[1]{ \subsection {#1} }
}

% Indonesian aliases used by the independent translated reader.
\newtheorem{Teorema}[fakt]{Teorema}
\renewcommand{\proofname}{Bukti}
